\documentclass[11pt]{amsart}
\usepackage[T1]{fontenc}
\usepackage[utf8]{inputenc}
\usepackage{lmodern,microtype,amsmath,amssymb,mathtools,mathrsfs}
\usepackage[margin=1.08in]{geometry}
\usepackage{enumitem,longtable,booktabs,array}
\usepackage{xcolor}
\usepackage[unicode,colorlinks=true,linkcolor=blue!50!black,citecolor=blue!50!black,urlcolor=blue!50!black]{hyperref}
\usepackage{url}
\urlstyle{same}
\Urlmuskip=0mu plus 1mu
\makeatletter
\g@addto@macro\UrlBreaks{\do\-}
\makeatother
\newtheorem{theorem}{Theorem}[section]
\newtheorem{proposition}[theorem]{Proposition}
\newtheorem{lemma}[theorem]{Lemma}
\newtheorem{corollary}[theorem]{Corollary}
\theoremstyle{definition}
\newtheorem{definition}[theorem]{Definition}
\theoremstyle{remark}
\newtheorem{remark}[theorem]{Remark}
\newcommand{\C}{\mathbb C}
\newcommand{\Q}{\mathbb Q}
\newcommand{\R}{\mathbb R}
\newcommand{\Z}{\mathbb Z}
\newcommand{\OO}{\mathcal O}
\newcommand{\PP}{\mathbb P}
\newcommand{\PE}{\overline{\operatorname{Eff}}}
\newcommand{\ddc}{i\partial\bar\partial}
\DeclareMathOperator{\ord}{ord}
\DeclareMathOperator{\rank}{rank}
\DeclareMathOperator{\Pic}{Pic}
\DeclareMathOperator{\Supp}{Supp}
\DeclareMathOperator{\Sym}{Sym}
\DeclareMathOperator{\Spec}{Spec}
\DeclareMathOperator{\divisor}{div}
\DeclareMathOperator{\trdeg}{trdeg}
\DeclareMathOperator{\im}{im}
\DeclareMathOperator{\codim}{codim}
\title[Cohomological transfer and anticanonical sections]{Cohomological transfer and equivariant anticanonical sections}
\author{OpenAI}
\date{September 26, 2026}
\hypersetup{pdfauthor={OpenAI},pdftitle={Cohomological transfer and equivariant anticanonical sections}}
\begin{document}
\begin{abstract}
Let $X$ be a smooth projective complex variety with smoothly semipositive anticanonical bundle. For any torus linearization of that bundle, we show that invariant sections in unbounded degrees with one fixed negative pseudoeffective error yield an invariant section in a positive untwisted degree. Combining the natural invariant index with compact-monodromy structure, we deduce that every such $X$ has a nonzero section of some positive anticanonical power.
\end{abstract}
\maketitle
\tableofcontents
\section{Introduction}\label{sec:introduction}
Let $X$ be a smooth connected projective variety over $\C$ and put $L=-K_X=\det T_X$. Anticanonical nonvanishing asks whether $H^0(X,mL)\ne0$ for some integer $m>0$. We assume that $L$ is smoothly semipositive: it has a smooth Hermitian metric with nonnegative Chern curvature. Cohomological constructions can produce sections with a fixed twist, which must then be removed. In the presence of a group action there is a second obstruction: the resulting section must have the required character. Our principal theorem addresses these two obstructions together.

Let an algebraic torus $T$ act on $X$. A $T$-linearization of a line bundle is a lift of this action to its total space, linear on each fiber. The tangent action gives the natural linearization of $L$, but tensoring that action by a character of $T$ gives another one. We write $P$ for $L$ equipped with any such linearization, and use additive notation for tensor products. Thus $H^0(X,mP)^T$ specifies both a degree and a character of an anticanonical section.

The input to our principal theorem is a sequence of sections of $M+iP$ with a fixed error $M$. The assumption on the error is that $-M$ is pseudoeffective: its numerical class lies in the closure of the effective cone.

\begin{theorem}[Invariant transfer]\label{inj:transfer}
Let $X$ be smooth connected projective, let $L=-K_X$ carry a smooth
semipositive metric, and let an algebraic torus $T$ act on $X$.
Let $P,M$ be linearized line bundles such that $P=L$ as ordinary
line bundles and $-M$ is pseudoeffective.  The linearization of $P$
may differ from the natural one on $L$.  Suppose that for an
unbounded $I\subset\mathbb Z_{>0}$ there are nonzero invariant
sections
\[
 s_i\in H^0(X,M+iP)^T\qquad(i\in I).
\]
Then $H^0(X,mP)^T\ne0$ for some integer $m>0$.
\end{theorem}


The theorem removes the fixed error without losing the specified character. It assumes neither vanishing of $H^1(X,\mathcal O_X)$ nor a condition on the Euler characteristic. Taking the trivial torus gives ordinary fixed-error conversion. With the natural linearization, invariant cohomology supplies an important source of the required twisted sections. In Section~\ref{sec:cohomological-input} we deduce that
\[
 \chi(X,\mathcal O_X)\ne0
 \quad\Longrightarrow\quad
 H^0(X,-mK_X)^T\ne0\text{ for some }m>0.
\]
The integer may depend on $X$. The distinction between natural and shifted linearizations remains essential in this deduction: the invariant anticanonical index is applied only with the natural one.

Combining this natural invariant nonvanishing with the compact-monodromy structure theorem of \cite[Theorem \textnormal{(Finite cover with compact torus monodromy)}]{CompanionH} removes the Euler-characteristic premise from the ordinary conclusion. The structure theorem separates the universal cover into a Euclidean factor, compact factors with parallel canonical frames, and a compact factor with no positive-degree holomorphic forms. After a finite cover, the remaining action on the last factor lies in a compact torus. Its invariant section therefore descends, and a finite norm returns a section to $X$. Section~\ref{sec:descent} proves the following consequence and constructs two additional descent maps, for cohomology and for twisted forms.

\begin{corollary}\label{thm:ordinary}
For every smooth connected projective complex $X$ with smoothly semipositive $-K_X$, some positive multiple $-mK_X$ has a nonzero section.
\end{corollary}

\subsection{Why a fibration removes the error}
Choose nonzero invariant input sections $s_i$. Ratios of products having the same linearized degree are invariant rational functions. Their common field determines a rational map from $X$ to a projective base. Along its generic fiber, the divisors of the $s_i$ vary affinely with $i$. Unboundedness forces their slope to be effective, so a quotient $s_j/s_i$ supplies a relative anticanonical section.

The main difficulty is to extend this relative section across the boundary. We work with a birational morphism $\mu:W\to X$ and a normal equidimensional model $f_W:W\to Y$, and put $P_W=\mu^*P$. Every vertical prime of $W$ then lies over a prime of $Y$. Subtracting the smallest normalized order over each base prime gives a regular section of $kP_W-f_W^*(kB)$, where $B$ is a rational divisor on $Y$ and $kB$ is integral. On a smooth resolution $\pi:V\to X$ through $W$, put $E=K_{V/X}$. The same ratio field shows that the invariant sections of $E+jk\pi^*P$ on the generic fiber of $V\to Y$ form a one-dimensional space over $\C(Y)$ for every $j\ge0$, including zero.

Two metrics now play different roles. Smooth adjoint fiber integrals in increasing degrees converge, after taking roots, to a fiberwise maximum. This gives a bounded semipositive metric on $kB$. At degree zero, pseudoeffectivity of $-M$ supplies a different metric with positive curvature in the base directions. Its fiber integral extends to $-K_Y+D_*$, where $D_*\ge0$ pulls back to an exceptional divisor over $X$. Its multiplier ideal $\mathcal J$ can be nontrivial. What matters is the inclusion $\mathcal O_Y(-D_*)\subseteq\mathcal J$ and the fact that tensoring by the bounded metric on $jkB$ leaves $\mathcal J$ unchanged.

Nadel vanishing therefore turns the section dimension of
$\mathcal O_Y(D_*+jkB)\otimes\mathcal J$ into a polynomial in $j$, including at zero. Its value there is positive, so it cannot vanish at every positive integer. A resulting section lifts by the relative section, and the exceptional correction disappears on $X$. Section~\ref{sec:models} prepares the model and the maximum metric; Section~\ref{inj:character-transfer} gives the complete transfer proof.

\subsection{A separate harmonic mechanism}
There is another way to detect the zero exponent, using restriction of cohomology to fibers. Its hypotheses and purpose differ from character-sensitive transfer. Let $f:V\to Y$ be a surjective morphism between smooth projective varieties, let $D_V$ have a smooth semipositive metric of curvature $\theta$, and let $Q$ have a continuous semipositive metric on $Y$. Suppose that on a nonempty ordinary open set where $f$ is a submersion, every null direction of $\theta$ is vertical. We prove in Theorem~\ref{har:thm:kernel} that
\[
 H^p\bigl(Y,R^qf_*(K_V+D_V)\otimes Q^k\bigr)=0
 \qquad(p>0,\ q,k\ge0).
\]
The proof approximates only the continuous metric pulled back from $Y$. A weak limit of harmonic representatives has all base canonical covectors as wedge factors, and its relative part detects a nonzero cohomology class on a general fiber. We apply decomposition to the smooth coefficient $D_V$ and then tensor the resulting derived isomorphism by $Q^k$. Thus the argument uses no decomposition theorem for a merely continuous coefficient upstairs.

We work through the natural-linearization application completely, then explain alternative constructions at the point where their geometric input changes. A fixed invariant system, null-pole functions, and determinant-semigroup ratios give different ways to choose the base. An ordinary curvature-null field leads instead to a reduced flat alteration, a line of multiplication maps, and a norm back to the original function field. Local strictness with an exact direct-image ideal and stabilization of actual direct-image lines provide two further vanishing mechanisms. Each construction is carried through to its section conclusion; the common analytic conclusions are reused only after their hypotheses have been verified.

\subsection{Context and analytic predecessors}
Anticanonical nonvanishing is motivated by the questions recorded by Yau \cite[Problem~75]{Yau1994}. Our conclusions concern a positive power, whose exponent may vary with the variety. Smooth semipositivity provides curvature identities unavailable from numerical nefness alone. Demailly--Peternell--Schneider established the universal-cover splitting and finite-cover Albanese fibration for a Hermitian semipositive anticanonical bundle \cite[Structure Theorem]{DPS1996}; Campana--Demailly--Peternell identified the remaining compact factors as rationally connected \cite[Theorem~1.4]{CDP2015}. The finite cover is a fiber bundle, so descent must still account for monodromy. We use the precise compact-torus structure and finite \'etale norm in \cite{CompanionH}, and construct the cohomology-valued and form-valued descent maps here.

In the broader nef setting, Lazi\'c--Matsumura--Peternell--Tsakanikas--Xie proved numerical effectivity of $-(K_X+\Delta)$ for projective log canonical threefold pairs with nef $-(K_X+\Delta)$, assuming that $X$ is $\Q$-factorial or has rational singularities \cite[Theorem~A]{LMPTX2023}. Numerical effectivity means that the numerical class is represented by an effective $\R$-divisor; it does not by itself provide a section of the specified line bundle.

M\"uller proves nonvanishing for projective klt pairs with nef log anticanonical divisor whose restriction to a general fiber of the maximal rationally connected fibration is semiample, and for every projective klt threefold pair with nef log anticanonical divisor \cite[Theorem~A and Corollary~B]{Muller2025}. His equivariant theorem gives natural invariant sections for semiample log anticanonical divisors on projective sub-log canonical pairs under commutative linear algebraic groups \cite[Theorem~C]{Muller2025}. Here we work with smooth projective varieties and algebraic tori. Smooth semipositivity and a nonzero Euler characteristic give natural invariant sections, while Theorem~\ref{inj:transfer} allows shifted linearizations with fixed-error input. The compact-monodromy descent then gives ordinary nonvanishing without the Euler-characteristic assumption.

The independent compact-complex index theorem, with its actual value at zero, is supplied by \cite[Invariant anticanonical index theorem]{CompanionI}. The passage from cohomology to twisted forms uses smooth semipositive hard Lefschetz, due to Mourougane \cite[Theorem~2.6 and Serre duality]{Mourougane1999}; we use the formulation \cite[Corollary~2.1.2]{DPS2001}. The determinant construction follows Lazi\'c--Peternell \cite[Lemma~4.1]{LP2018} and its later use in \cite[Lemma~5.1]{LMPTX2023}; we prove the unbounded-twist argument in the form needed here, including the numerically trivial case. The inverse determinant is pseudoeffective by the cotangent-tensor theorem of Lazi\'c--Matsumura--Peternell--Tsakanikas--Xie \cite[Theorem~4.1]{LMPTX2023}, which develops Ou's generic-nefness argument \cite[Theorem~1.4]{Ou2023}.

Berndtsson's smooth direct-image theorem \cite{Berndtsson2009} gives positivity of the fiberwise adjoint metrics. We obtain the continuous maximum metric by taking their normalized moment limits, and derive local horizontal strictness from his curvature formula. For singular coefficients, Berndtsson--P\u aun proved direct-image positivity over smooth families \cite[Theorem~3.5]{BP2008}; we use P\u aun--Takayama's multiplier-ideal formulation and extension to the actual torsion-free direct image \cite[Theorems~3.3.4--3.3.5]{PT2018}. Positivity of an invariant line is obtained from a specified holomorphic quotient; it is never inferred for an arbitrary singular subbundle. Fujino--Matsumura injectivity \cite{FM2016} and Nadel vanishing \cite{Nadel1990} supply the cohomological ending. Guan--Zhou strong openness \cite{GuanZhou2015} is used to keep the relevant ideal fixed when a strict metric is mixed with another one.

The harmonic argument belongs to a different regularity regime. Richberg smoothing \cite{Richberg1968} provides uniform approximations of continuous base weights, and Fujisawa's derived decomposition \cite[Theorem~1]{Fujisawa2015}, obtained from Takegoshi's harmonic representatives, applies to the smooth semipositive coefficient on $V$. Related vanishing theorems impose global curvature domination: Fujino \cite[Theorem~1.3]{Fujino2018} uses a smooth semipositive metric on the pullback of an ample base line, whereas Matsumura \cite[Theorem~1.3]{Matsumura2018} uses the pullback of a base K\"ahler form. Our harmonic proof starts instead from the kernel inclusion on a nonempty ordinary open set and treats the continuous base twist by uniform approximation. The reduced-flat construction uses weak semistable reduction \cite{AK2000}; its finite-group descent and normal-target norm are kept distinct from a finite-flat norm-line construction.

\subsection{Conventions and reading path}
We write $\sim$, $\sim_{\Q}$, and $\equiv$ for linear, rational linear, and numerical equivalence. A local frame of squared norm $e^{-\phi}$ has curvature $i\partial\bar\partial\phi$; factors of $2\pi$ are included when viewing curvature as a first Chern form. A continuous semipositive metric has continuous plurisubharmonic weights. Locally bounded weights have both upper and lower bounds and need not be continuous.

The action on sections is $(g\cdot s)(gx)=g(s(x))$. Algebraic-torus invariants agree with invariants under its maximal compact torus. Averaging smooth metric weights and K\"ahler forms over the compact torus preserves semipositivity. A rank assertion on a dense base-change open is not an assertion on singular fibers.

The shortest route to Corollary~\ref{thm:ordinary} is the transfer proof in Sections~\ref{sec:models}--\ref{inj:character-transfer}, its natural invariant application in Section~\ref{sec:cohomological-input}, and invariant-section descent in Section~\ref{sec:descent}. Section~\ref{inj:extremal-weights} also obtains exact source characters. The remaining constructions explain different ways to convert cohomology into sections. Section~\ref{inj:injectivity} gives an independent ordinary proof through singular injectivity; Sections~\ref{har:sec:harmonic}--\ref{har:sec:alteration} develop harmonic restriction and its geometric inputs; Sections~\ref{inj:localized} and~\ref{inj:stabilized} use exact multiplier ideals and actual direct-image lines. The final section also gives the stronger tensor-holonomy calculation.

The table compares the inputs and the step that completes each conversion. The source-character constructions in Sections~\ref{inj:extremal-weights} and~\ref{har:sec:extremal} additionally keep track of an extremal character.
\begin{center}
\small
\begin{tabular}{@{}>{\raggedright\arraybackslash}p{.25\textwidth}>{\raggedright\arraybackslash}p{.29\textwidth}>{\raggedright\arraybackslash}p{.37\textwidth}@{}}
\toprule
Input & Base construction & Conversion \\
\midrule
Unbounded shifted invariant sections; $-M$ pseudoeffective & Common section-ratio field & Fixed-ideal vanishing; Theorem~\ref{inj:transfer} \\
Natural invariant cohomology; nonzero Euler value & Fixed invariant system & Harmonic restriction; Section~\ref{har:sec:eulerroute} \\
Ordinary cohomology in one degree at unbounded powers & Curvature-null field and reduced flat alteration & Descent of a continuous line and its injection; Theorem~\ref{thm:alteration-conversion} \\
No positive structure-sheaf cohomology & Null-pole field or determinant semigroup & Maximum metric and harmonic restriction; Section~\ref{har:sec:fields} \\
No positive structure-sheaf cohomology & Null-pole field, expressed by degree zero & Local strictness for the exact direct-image ideal; Section~\ref{inj:localized} \\
Unbounded natural invariant twisted sections; $-M$ pseudoeffective & Toroidal invariant Iitaka model & Affine boundary orders and strong openness; Section~\ref{inj:stabilized} \\
\bottomrule
\end{tabular}
\end{center}

\section{Preparing a fibration and measuring its relative section}\label{sec:models}

A section regular on the generic fiber may still have poles over the boundary. The following model separates the horizontal and vertical tests for regularity; a smooth resolution will be used only for fiber integrals.

\begin{lemma}[Equidimensional models]\label{val:flat-model}
Let $X\dashrightarrow Y$ be a dominant rational map of integral projective
varieties over $\mathbb C$.  There are a smooth projective birational model
$Y'$ of $Y$ and an integral projective flat family $X_0\to Y'$ with a
birational morphism to the graph of the original map.  Every irreducible
component of every fiber has dimension $\dim X-\dim Y$.  The normalization
$V$ of $X_0$ is equidimensional over $Y'$; in particular, every vertical
prime divisor of $V$ maps onto a prime divisor of $Y'$.  If
$\mathbb C(Y)$ is algebraically closed in $\mathbb C(X)$, the morphism
from $V$, and from any smooth resolution of $V$, to $Y'$ has connected
fibers.  If a torus acts on $X$ and the rational map is invariant, these
constructions and a projective resolution of $V$ can be made equivariant,
with trivial action on the base.
\end{lemma}

\begin{proof}
Embed the projective graph over $Y$ in $\mathbb P^N\times Y$.
Over the flat open set the fibers define a map to the projective
Hilbert scheme \cite[Theorem~3.2 and the Hilbert-scheme specialization
after Proposition~3.8]{Grothendieck1961Hilbert}.
Resolve the closure of its graph, also resolving the base, and pull back the
universal family.  Flatness over an integral base embeds each affine
coordinate ring in its generic localization.  The generic fiber is
integral, so these coordinate rings have neither nilpotents nor zero
divisors.  Thus the pulled-back family is integral and is the closure of
the original family over the flat open set.  Its map to the graph is
birational.

Constancy of the Hilbert polynomial bounds the dimension of every fiber
above by the generic relative dimension.  At any fiber component the
height bound for the ideal generated by local parameters of the smooth
base gives the opposite bound.  This proves purity of fiber dimension.
Normalization is finite, so fiber dimensions and dimensions of images of
components are preserved.  A prime divisor mapping into codimension at
least two in $Y'$ would have dimension at most
$\dim Y'-2+(\dim X-\dim Y')=\dim X-2$, a contradiction.

For connectedness apply Stein factorization to the normal total space.
Its finite intermediate base has function field the algebraic closure
of $\mathbb C(Y')$ in $\mathbb C(X)$; the asserted field condition and
normality of $Y'$ make the finite birational morphism an isomorphism.
The same reasoning applies after resolution.  In the equivariant case,
the family over the original open set is stable under the action.
Its closure and normalization are therefore stable, and equivariant
resolution in characteristic zero supplies the last assertion.
\end{proof}


\subsection{Rank-one adjoints and a continuous maximum metric}\label{har:sec:maximum}

Our first task is to turn a rational relative section into a metrized line on the base. This construction will serve both the character-sensitive transfer and the harmonic applications.
Its normal equidimensional model is used for divisorial normalization
and continuity. A resolution supplies smooth coefficients for the
direct-image theorem; that resolution need not be equidimensional.

\begin{proposition}[Divisorial normalization and maximum metric]\label{har:prop:maximum}
Let \(X,Y\) be smooth connected projective, let \(L=-K_X\) have a smooth
semipositive metric, and suppose
\[
 V\xrightarrow{\nu}W\xrightarrow{\mu}X,\qquad g:W\to Y,
 \qquad f=g\nu,
\]
where \(W\) is normal, \(g\) is projective surjective equidimensional,
\(\mu\) is birational, and \(V\) is a smooth projective resolution.
Set \(D_V=(\mu\nu)^*L\) and \(E=K_{V/X}\), and let \(s_E\) be the canonical Jacobian section of \(E\).
Suppose a nonzero rational section \(s\) of \(b\mu^*L\), \(b>0\), has
no horizontal poles and, on the generic fiber $V_\eta$, satisfies
\[
 H^0(V_\eta,E+kbD_V)=\C(Y)\,s_Es^k
 \qquad(k\ge0).
\]
Alternatively, suppose a torus acts compatibly on the diagram, trivially
on $Y$, and let $P$ be $L$ with any chosen linearization. Put
$P_V=(\mu\nu)^*P$. In this case it suffices that $s$ be an invariant
rational section of $b\mu^*P$ and that the displayed condition be
replaced by
\[
 H^0(V_\eta,E+kbP_V)^T=\C(Y)\,s_Es^k
 \qquad(k\ge0),
\]
where $E$ carries its natural Jacobian action.

After replacing \(b,s\) by a multiple and a power, there are a
line \(Q\) on \(Y\) and a regular injection
\[
 g^*Q\longrightarrow b\mu^*L
\]
whose squared fiber-maximum norms define a positive continuous metric
with plurisubharmonic weights on \(Q\).
In the invariant case $Q$ has trivial linearization and the injection
$g^*Q\to b\mu^*P$ is equivariant.
\end{proposition}

\begin{proof}
In the invariant case average a smooth semipositive metric on $L$ over
the maximal compact torus. The chosen linearization of $P$ differs
from the natural one by a character, which is unitary on this compact
torus, so the averaged metric is invariant for both actions.

The rank assumption at $k=0$ forces the generic fiber to be geometrically connected. Indeed its constant functions, multiplied by $s_E$, inject into that adjoint space. In the invariant case the connected torus acts trivially on the finite Stein base, so all these constants are invariant. Their dimension is therefore one. Smooth general fibers are consequently connected and irreducible.

For every prime divisor \(B\) of \(Y\), put
\[
 c_B=\min_{\Gamma\mapsto B}
       \frac{\operatorname{ord}_{\Gamma}(s)}
            {\operatorname{ord}_{\Gamma}(g^*B)}.
\]
Properness and equidimensionality supply divisorial lifts of \(B\);
all vertical prime divisors lie over base primes. Only finitely many
\(c_B\) are nonzero. Clear denominators and use the resulting integral
divisor as \(Q\). Sending its rational canonical section to the powered
\(s\) gives a regular bundle map by normality. Its effective Cartier
zero divisor \(G\) omits at least one component above every base prime.
In the invariant case all these divisors are invariant, the base line carries trivial linearization, and the map is to \(b\mu^*P\). In an eigensection version, one
instead twists the base linearization by the section character.

For a local frame \(e\) of \(Q\), denote its image by \(s_e\) and put
\begin{equation}\label{har:eq:max}
 M_e(y)=\max_{x\in W_y}|s_e(x)|^2.
\end{equation}
This is continuous. Properness gives upper semicontinuity. For lower
semicontinuity use that an equidimensional holomorphic map from a pure
dimensional normal space to a smooth base is open. To verify openness
at a point, cut the source by as many general local hyperplanes as the
fiber dimension, isolating the point in its fiber. The slice germ has
dimension at least that of the base and is finite over that smooth
base germ. Its image is therefore the entire base germ. This can be
done inside any prescribed source neighborhood.

The locus \(Z_0\) where \(M_e\) vanishes is algebraic closed: its
complement is the constructible image \(g(W\setminus\operatorname{Supp}G)\),
and continuity makes \(Z_0\) closed in the classical topology; a constructible
classically closed subset is Zariski closed. By the omitted-lift
property, \(Z_0\) has codimension at least two.

We show first that \(-\log M_e\) is plurisubharmonic on a dense open.
Work on a fixed smooth-fibration open for \(f\), where the canonical
exceptional section and \(s_e\) are not identically zero on any fiber
component. Take a local frame \(\eta\) of \(K_Y\). The section
\[
  (s_E/f^*\eta)s_e^k
 \quad\text{of }K_{V/Y}+(kb+1)D_V
\]
generates the rank-one adjoint space at the generic point. In the shifted equivariant case its coefficient is \(D_V+kbP_V\); as a metrized ordinary line it is still \((kb+1)D_V\). Berndtsson's
smooth direct-image theorem \cite[Theorem~1.2]{Berndtsson2009} applies on
the bundle and base-change locus with smooth semipositive coefficient
\((kb+1)D_V\). In the invariant case the invariant line is an
\emph{orthogonal quotient} of the full direct image: the compact torus
acts unitarily, and the sum of its nonzero-weight spaces is its
orthogonal complement. It therefore has semipositive curvature.
This quotient argument, rather than positivity of arbitrary sublines,
is what permits taking invariants.

For $k\ge1$ the resulting plurisubharmonic function is
\begin{equation}\label{har:eq:moments}
 -\frac1k\log\int_{V_y}|s_e|^{2k}\,d\mu_y,
\end{equation}
where \(\mu_y\) comes from \(s_E/f^*\eta\) with the metric of \(D_V\).
It is smooth, nonnegative, of full support, and independent of \(k\).
Each integral is smooth positive on our fixed open, so the curvature
inequality extends there from its possibly \(k\)-dependent base-change
open. Normalized to probability measures, its \(1/k\)-powers increase
to \(M_e\). They converge locally uniformly: a point just below the
positive fiber maximum has a neighborhood of uniformly positive
\(\mu_y\)-mass on nearby fibers; compactness supplies finitely many such
neighborhoods, while bounded total mass gives the upper estimate.
Thus \eqref{har:eq:moments} converges locally uniformly to \(-\log M_e\).

By continuity this weight extends plurisubharmonically across the bad
analytic loci wherever \(M_e>0\). It cannot tend to \(+\infty\) across
\(Z_0\) of codimension at least two. Through a point of $Z_0$ choose
a small analytic disc whose boundary misses $Z_0$. Nearby parallel
discs have their boundaries in a fixed compact set off $Z_0$.
Since $Z_0$ has codimension at least two, a dense set of these discs
misses it entirely. The maximum principle therefore bounds the weight
uniformly at their centers, and continuity off $Z_0$ supplies the
same bound there. A zero of the continuous \(M_e\) would contradict
this bound. Hence \(M_e>0\) everywhere, and the weight is continuous
plurisubharmonic. Its transformation rule under a change of frame is
exactly the metric transformation rule for \(Q\).
\end{proof}

\section{Character-sensitive transfer with a fixed multiplier ideal}
\label{inj:character-transfer}

The next conversion allows a character shift of the anticanonical
linearization.  Its proof uses a common field of invariant section
ratios and a singular metric with positive curvature in the base
directions.  The boundary correction on the base has exceptional
pullback, and its multiplier ideal stays fixed while the exponent
varies.

\begin{proof}[Proof of Theorem~\ref{inj:transfer}]
\smallskip\noindent\textit{The common ratio field.} Put $n=\dim X$ and write $N_i=M+iP$.  For each $i$, ratios of invariant sections of
the same positive multiple of $N_i$ form a field $Q_i$.
For $p<q<b$ in $I$, the identity
\[
 (b-p)N_q=(b-q)N_p+(q-p)N_b
\]
holds with linearizations.  Multiplication by suitable powers of
the fixed $s_p,s_b$ puts any ratio for $N_p$ or $N_b$ into $Q_q$.
This retains the ratio itself: for a ratio in degree $aN_p$, multiply
its numerator and denominator by
$s_p^{a(b-q-1)}s_b^{a(q-p)}$, and use the analogous formula for
a ratio in $aN_b$.  Taking a further index beyond any prescribed
two interior indices proves that all $Q_i$ with $i>\min I$ coincide;
denote this field by $Q$.

Fix the first two indices $p<q$.  For $b>q$ the ratio
\begin{equation}\label{inj:ratio-field-identity}
 \frac{s_q^{b-p}}{s_p^{b-q}s_b^{q-p}}
\end{equation}
belongs to $Q$.  If $\operatorname{trdeg}_{\mathbb C}Q=0$,
it is constant.  The section divisors are affine in $i$, and their
slope is effective because the indices are unbounded.  Thus
$s_q/s_p$ is a regular invariant section of $(q-p)P$.

Otherwise put $e=\operatorname{trdeg}Q>0$.  Choose $e$ independent
ratios to define a map to $\mathbb P^e$, resolve its graph and take
Stein factorization. The Stein field is the relative algebraic closure
in $\mathbb C(X)$ of the field of the chosen ratios, a finite extension
of that field. It contains $Q$, since $Q$ is algebraic over the chosen
ratio field, and is relatively algebraically closed in $\mathbb C(X)$.
The connected torus acts
trivially on the Stein base because its orbits lie in finite fibers
over a base with trivial action.  As before take a smooth projective
base $Y$, a normal equidimensional $W\to Y$ with a birational
morphism $W\to X$, and an equivariant smooth resolution $V\to W$,
with maps $\mu:W\to X$, $f_W:W\to Y$, and total maps $\pi:V\to X$ and $f:V\to Y$. Put $P_W=\mu^*P$.  The map $f$ has
connected fibers by the Stein construction and normality of $Y$.
Write $L'=\pi^*L$ and
$P'=\pi^*P$.

For each interior index $i$ and every divisor $H'$ on $Y$,
\begin{equation}\label{inj:ratio-domination}
 H^0(V,a\pi^*N_i-f^*H')^T\ne0
 \quad\text{for some integer }a>0.
\end{equation}
Indeed the chosen independent ratios belong to $Q_i$.  Put their
numerators and denominator in one common degree by products of
invariant sections.  On the resolved graph they have a common
invariant factor whose line is that multiple of $N_i$ minus the
pullback of $\mathcal O_{\mathbb P^e}(1)$.  The latter line is
big on $Y$, so a sufficiently large multiple dominates $H'$.

\smallskip\noindent\textit{A regular relative section.} We next remove the vertical poles while preserving the character. Let $D_i$ be the section divisor on $W$.  Since
\eqref{inj:ratio-field-identity} comes from $Y$, its horizontal
divisor is zero.  The affine relations show that the horizontal
part of $(D_q-D_p)/(q-p)$ is effective.  Normalize the vertical
part by subtracting the minimum coefficient divided by pullback
multiplicity above each base prime, and let $B$ be the resulting
rational base divisor.  There is an integer $k>0$, divisible by
$q-p$ and clearing the coefficients of $B$, for which
\begin{equation}\label{inj:gamma}
 \gamma=(s_q/s_p)^{k/(q-p)}(f_W^*1_{kB})^{-1}
 \in H^0(W,kP_W-f_W^*(kB))^T
\end{equation}
is regular.  Normality checks regularity in codimension one, and
equidimensionality accounts for every vertical prime.  Above every
base prime $\gamma$ is nonzero at the general point of at least
one component.  We also write $\gamma$ for its pullback to $V$.

\smallskip\noindent\textit{All adjoint ranks, including zero.} The ratio field must also control the fiber integrals. Let $E=K_{V/X}\ge0$.  Its Jacobian section $1_E$ is invariant
for the natural linearization on $L'$, irrespective of the character
shift in $P'$.  For all $j\ge0$,
\begin{equation}\label{inj:transfer-rank}
 \operatorname{rank}
 \bigl(f_*\mathcal O_V(E+jkP')\bigr)^T=1.
\end{equation}
The section $1_E\gamma^j$ gives the lower bound over the generic
fiber.  If the rank were larger, globalize two invariant sections
$u_1,u_2$ independent over $\mathbb C(Y)$ after a twist $f^*(aH)$,
where $H$ is ample.  Choose $i,b\in I$ with $b>i+jk$.  Multiply
$u_1^k$ and $u_1^{k-1}u_2$ by
$s_i^k\gamma^{b-i-jk}$.  Both results are sections of
\[
 kE+k\pi^*N_b+
 f^*\bigl(kaH-(b-i-jk)kB\bigr),
\]
with ratio $u_2/u_1$.  Apply \eqref{inj:ratio-domination} at
index $b$ to cancel this base twist by the same invariant multiplier
in the numerator and denominator.  Removing the exceptional twist
$kE$ and pushing to $X$ leaves two invariant sections of an equal
multiple of $N_b$.  Normality extends across the image of the
exceptional locus.  Their ratio must be in $Q\subset\mathbb C(Y)$,
contrary to independence.  This proves \eqref{inj:transfer-rank}.

\smallskip\noindent\textit{The bounded metric and the horizontally positive metric.} We now construct the two metrics needed for interpolation. Average a smooth semipositive metric $h_0$ on $L$ over the maximal
compact torus.  It is also invariant for $P$, because the two
linearizations differ by a character, unitary on that compact torus.
Equation \eqref{inj:ratio-domination} and pseudoeffectivity of $-M$
give $L'-\eta f^*H$ pseudoeffective for some $\eta>0$ and ample
$H$.  Taking a positive current, averaging, and mixing a small
positive fraction of its potential with $h_0$ gives an invariant
singular metric $h$ on $L'$ with
\begin{equation}\label{inj:strict-integrable}
 i\Theta_h\ge\delta f^*\omega_Y,
 \qquad \mathcal I(h)=\mathcal O_V,
 \qquad\delta>0.
\end{equation}
Skoda integrability and compactness give the last equality.  The
metric is bounded below by a positive constant times $h_0$ on
compact coordinate patches: its potential relative to the smooth
metric is locally bounded above.

Apply Proposition~\ref{har:prop:maximum} with the normal model $W$, the chosen linearization $P$, and the relative section $\gamma$ of $kP_W-f_W^*(kB)$. Equivalently use the rational section $(s_q/s_p)^{k/(q-p)}$ before normalization. Equation~\eqref{inj:transfer-rank} gives every invariant adjoint rank, the Jacobian has its natural invariant action, and the averaged smooth metric is invariant also for $P$. The vertical minima defining $kB$ are exactly those in that proposition. Its smooth coefficient is $L'+jkP'$, ordinary $(1+jk)L'$, so the orthogonal quotient argument applies without any singular-subbundle claim. We obtain a continuous, in particular locally bounded, semipositive metric on $kB$.

\smallskip\noindent\textit{A fixed ideal on the base.} The continuous base metric just obtained controls the varying power. The singular degree-zero metric will supply the positive constant term.
For $j=0$ use the singular coefficient metric $h$.
The singular adjoint direct-image theorem
\cite[Setup 3.2.1 and Theorem 3.3.4]{PT2018} applies over the
smooth, locally free base-change locus.  Its generic multiplier-ideal
inclusion is an isomorphism by \eqref{inj:strict-integrable}; fiber
integrability is an almost-everywhere assertion by Fubini.
The invariant summand is an orthogonal holomorphic quotient for
the compact-invariant metric.  Thus its line metric has
semipositive curvature.  Subtract a local potential for
$\delta\omega_Y$ from the coefficient weight to see that its
curvature is in fact at least $\delta\omega_Y$.

In a frame $\xi$ of $-K_Y$, denote the squared norm of
$1_Ef^*\xi$ by $F_\xi(y)$.  With base volume in the dual frame,
this is the pushforward density of the integrable absolute adjoint
form $1_E$ with metric $h$.  Hence $F_\xi$ is locally integrable
on the whole base chart.  We next extend its negative logarithm
with an explicit exceptional correction.

\smallskip\noindent\textit{Boundary extension with exceptional support.} For a prime boundary divisor $P$ choose a prime $Q\subset V$
mapping onto it, of pullback multiplicity $l$ and Jacobian order
$a=\operatorname{ord}_Q(E)$.  At general points write
$f=(z_1^l,z_2,\ldots,z_e)$.  Integrating on one coordinate patch,
using $h\ge c h_0$, gives almost everywhere
\begin{equation}\label{inj:density-bound}
 F_\xi(y)\ge c'|y_1|^{2((a+1)/l-1)}.
\end{equation}
The exponent includes the transverse Jacobian factor
$|z_1|^{-2(l-1)}$. If a nonexceptional component dominates $P$, choose
it for the estimate. Then $a=0$ and the coefficient of $P$ in the
correction is zero. Otherwise every component dominating $P$ is
exceptional; choose a nonnegative integer coefficient at least
$(a+1)/l-1$. Summing
over the finitely many boundary primes gives $D_*\ge0$ with
$f^*D_*$ exceptional over $X$.

For the last assertion one must also check primes mapping to
codimension at least two on $Y$.  The equidimensional normal model
$W$ has relative dimension $n-e$, so for $Z\subset Y$ of
codimension at least two,
\[
 \dim f_W^{-1}(Z)\le\dim Z+n-e\le n-2.
\]
Any prime of $V$ over $Z$ therefore maps to codimension at least
two on $W$ and on $X$. Every component dominating a base prime in
$\Supp D_*$ is exceptional by the rule defining that support. These two checks
account for every component of $f^*D_*$.

On $-K_Y+D_*$ define the metric by
$|\xi1_{D_*}|^2=F_\xi$ on the dense open.
Its ordinary-frame weight is
$-\log F_\xi+\log|g_{D_*}|^2$.  Estimate
\eqref{inj:density-bound} bounds this weight above at general
boundary points. Subtract a local smooth potential for
$\delta\omega_Y$. The resulting plurisubharmonic weight extends
across each boundary divisor away from a codimension-two analytic
subset. Hartogs extension for plurisubharmonic functions then fills
that subset. Restoring the smooth potential gives a metric on all of
$Y$ with curvature at least $\delta\omega_Y$.
If $\mathcal J$ is its multiplier ideal, then
\begin{equation}\label{inj:fixed-ideal-inclusion}
 \mathcal O_Y(-D_*)\subseteq\mathcal J,
\end{equation}
since multiplying the metric coefficient by $|g_{D_*}|^2$
recovers the locally integrable density $F_\xi$.

\smallskip\noindent\textit{Interpolation and return to $X$.} Tensor by the bounded semipositive metric on $jkB$.  For every
$j\ge0$ the multiplier ideal remains the same $\mathcal J$,
which is not assumed trivial.  Nadel vanishing yields
\[
 H^v\bigl(Y,\mathcal O_Y(D_*+jkB)\otimes\mathcal J\bigr)=0
 \qquad(v>0,\ j\ge0).
\]
Put $\mathcal F=\mathcal O_Y(D_*)\otimes\mathcal J$.
Its Euler characteristic after twisting by $jkB$ is a polynomial in
$j$. The inclusion \eqref{inj:fixed-ideal-inclusion} gives
$\mathcal O_Y\hookrightarrow\mathcal F$, so at zero this polynomial
equals $h^0(Y,\mathcal F)>0$. A nonzero section therefore exists for
some $j>0$. Pull it back
and multiply by $\gamma^j$ to obtain an invariant section of
$jkP'+f^*D_*$.  Removing the exceptional twist and pushing down
gives the required section of $jkP$.
\end{proof}

\section{Natural linearization and cohomological input}\label{sec:cohomological-input}
Write $n=\dim X$ and $\chi^T(X,\mathcal O_X)=\sum_q(-1)^q\dim H^q(X,\mathcal O_X)^T$. The transfer theorem permits a chosen character. To obtain its input from cohomology, we now use the natural anticanonical character. The following result is the exact input from the index companion \cite[Invariant anticanonical index theorem]{CompanionI}.

\begin{proposition}[Invariant anticanonical index]\label{har:prop:index}
Let $F$ be a compact complex manifold with a holomorphic action of a compact real torus $T$. Give $D=-K_F$ its natural linearization. There are an integer $a>0$ and a polynomial $P\in\Q[t]$ such that
\[
 \sum_q(-1)^q\dim H^q(F,maD)^T=P(m)\quad(m>0),\qquad
 P(0)=\sum_q(-1)^q\dim H^q(F,\mathcal O_F)^T.
\]
The statement includes the trivial torus, an empty fixed locus, and zero-dimensional fixed components; it assumes no positivity, projectivity, or K\"ahler condition.
\end{proposition}

In particular the value at zero is the actual index, not a formal extrapolation of an eventual polynomial. We use this theorem to select a fixed cohomological degree on an unbounded sequence.

\begin{lemma}[Forms and determinant sections]\label{har:lem:forms}
Let \(X\) be smooth connected projective, let \(L=-K_X\) have a smooth
semipositive metric, and let an algebraic torus \(T\) act on \(X\).
If \(\chi^T(X,\mathcal O_X)\ne0\), then for some fixed \(q\) there are
nonzero invariant sections of
\(\Omega_X^{n-q}\otimes L^{m+1}\) for unbounded positive integers \(m\).
If these do not already give a nonzero invariant section of a positive power of $L$, there
is an integer $e>0$, a linearized line bundle \(M\), a nonzero map
\(M\to\bigwedge^e\Omega_X^{n-q}\), and invariant sections
\[
 0\ne s_i\in H^0(X,M+a_iL)^T,\qquad a_i\longrightarrow\infty.
\]
\end{lemma}

\begin{proof}
The polynomial in Proposition~\ref{har:prop:index} is not zero. A fixed
cohomological degree therefore has nonzero invariant cohomology along an
unbounded sequence. Average a Kähler form and the metric weights over
the maximal compact torus, and denote the resulting curvature of \(L\)
by \(\theta\). The smooth semipositive hard Lefschetz theorem
\cite[Corollary~2.1.2]{DPS2001}, applied to \(mL=K_X+(m+1)L\),
gives the asserted forms. Concretely, if \(u\) is harmonic
of type \((n,q)\) with values in \((m+1)L\), Bochner--Kodaira gives
\[
 0=\|(D')^*u\|^2+\langle[(m+1)\theta,\Lambda]u,u\rangle.
\]
Both terms are nonnegative. Applying the complex-linear star on form
indices to $u$ therefore gives a holomorphic form of type \((n-q,0)\).
The construction is equivariant.
For \(q=0\) it already yields invariant sections of \(mL\); for \(q=n\)
the holomorphic form itself is a section of \((m+1)L\). If
\(r=\max\operatorname{rank}\theta=n\) and
\(q>0\), the curvature term makes this form vanish on a nonempty open set,
hence everywhere, so this case cannot occur.

The determinant extraction follows the method of
\cite[Lemma~5.1]{LMPTX2023} and \cite[Lemma~4.1]{LP2018}.
We spell out the unbounded-twist step, which here also applies when
$L$ is numerically trivial. Choose a rational frame of $L$ and divide
each form by the corresponding power of that frame. Their span over
$\C(X)$ is a subspace of the generic fiber of $\Omega_X^{n-q}$,
of some rank $e$. Changing the frame rescales each vector by a nonzero
rational function, so this subspace is intrinsic. Saturate its top
exterior power in $\bigwedge^e\Omega_X^{n-q}$ and call the resulting
rank-one subsheaf $M$. It is reflexive on smooth $X$, hence a line
bundle. The generic span is torus-stable, so its saturation carries the
induced linearization.

Choose finitely many of the original forms as a generic basis. Every
wedge of $e$ original forms lies generically in $M$ tensored by the
sum of their powers of $L$; saturation makes it a regular section of
that line. In the fixed basis, one coordinate of the unbounded family
is nonzero for unbounded $m$. Wedge those forms with the other $e-1$
fixed basis forms. These wedges are nonzero and invariant, and their
twist exponents are $m+1$ plus a fixed sum. They give the required
sections with $a_i\to\infty$.
\end{proof}


\begin{corollary}[Natural invariant nonvanishing]\label{thm:main}
Let $X$ be smooth connected projective over $\C$, with smoothly semipositive $L=-K_X$. If an algebraic torus $T$ acts on $X$ and $\chi(X,\mathcal O_X)\ne0$, then $H^0(X,mL)^T\ne0$ for some integer $m>0$, with the natural linearization. The torus may be trivial.
\end{corollary}
\begin{proof}
Every element of the connected torus is homotopic to the identity, so its action on singular cohomology is trivial. By Hodge decomposition its action on every $H^q(X,\mathcal O_X)$ is trivial. Consequently the invariant index at zero equals $\chi(X,\mathcal O_X)\ne0$.
Lemma~\ref{har:lem:forms} supplies either an invariant section directly or invariant sections of $M+a_iL$ with $a_i\to\infty$ and a nonzero map from $M$ into a positive tensor power of $\Omega_X^1$. The cotangent-tensor theorem \cite[Theorem~4.1]{LMPTX2023}, applied to the smooth zero-boundary pair with nef $-K_X$, says that $-M$ is pseudoeffective. Apply Theorem~\ref{inj:transfer} with $P=L$ naturally linearized.
\end{proof}

The harmonic constructions need additional information about the zero
divisors of these determinant sections. The following degree argument
provides it directly from the cotangent inclusion, independently of the
pseudoeffectivity theorem used in Corollary~\ref{thm:main}.

\begin{lemma}[Curvature-null determinant divisors]\label{har:lem:null-divisors}
Let $X$ be a smooth connected projective variety of dimension $n$, and
let $L=-K_X$ have smooth semipositive curvature $\theta$ of maximal
rank $r<n$. Suppose a line bundle $M$ admits a nonzero map
$M\to\bigwedge^e\Omega_X^p$ for integers $e,p>0$. Then
$c_1(M)G^{n-1}\le0$ for every Kähler class $G$. If
$0\ne s\in H^0(X,M+aL)$ for an integer $a$, its zero divisor $D$
satisfies
\[
 [D][\theta]^r[\omega]^{n-r-1}=0
\]
for every Kähler form $\omega$. Each component has zero intersection
separately, and $\theta^r$ vanishes on its smooth locus.
\end{lemma}

\begin{proof}
For every Kähler class \(G\), Yau's prescribed-Ricci theorem
\cite{Yau1978} gives a Kähler metric in \(G\) with Ricci form \(\theta\).
Kähler curvature symmetry identifies the mean curvature of \(T_X\)
with Ricci; the induced mean curvature of \(\bigwedge^e\Omega_X^p\)
is nonpositive. If \(\deg_G M>0\), solve a scalar Laplace
equation to give \(M\) a metric of positive constant mean curvature.
Its inverse tensor the cotangent construction then has negative mean
curvature, contradicting the Bochner formula for its nonzero section.
Thus \(c_1(M)G^{n-1}\le0\). Put \(G=[\theta]+t[\omega]\), divide by
\(t^{n-1-r}\), and let \(t\downarrow0\). Since \(\theta^{r+1}=0\),
\[
 c_1(M)[\theta]^r[\omega]^{n-1-r}\le0.
\]
The same intersection for \(M+aL\) is nonnegative by effectivity and
smooth semipositivity, and the contribution from $aL$ is zero. It is
consequently zero. Each component has zero intersection separately;
the nonnegative form $\theta^r\wedge\omega^{n-r-1}$ therefore vanishes
on its smooth locus, which gives the last assertion.
\end{proof}

\begin{remark}[The singular-norm proof of the degree sign]
The determinant degree calculation also has a useful current proof. In a K\"ahler metric with the prescribed semipositive Ricci form, let $e$ be the local image of a frame of $M$ in the cotangent construction. Its mean curvature is nonpositive. Chern's norm differentiation formula makes the trace of $i\partial\bar\partial\log(|e|^2+\varepsilon)$ nonnegative: the negative mean-curvature term has the required sign, and the derivative term is nonnegative by Cauchy--Schwarz. The logarithms converge locally in $L^1$, by comparison with logarithms of holomorphic coefficients. The resulting trace-positive current represents $-c_1(M)$. Integrating proves $c_1(M)G^{n-1}\le0$, including possible codimension-two zeros of the injected frame. This retains a singular-norm proof of the degree sign independently of the constant-mean-curvature contradiction used above.
\end{remark}

\section{Exact source characters from the transfer theorem}
\label{inj:extremal-weights}

A source component is the fixed component whose attracting set as $t\to0$ is open. The following calculation gives the exact character that will be transferred; it also supplies the weight input to the later harmonic method.

\begin{lemma}[The character at a source]\label{lem:source-character}
Let a torus $T$ act on a smooth connected projective $X$ with a linearized ample line. Let $U$ be a $T$-linearized line, let $\lambda$ be a generic integral cocharacter with $X^\lambda=X^T$, and let $S$ be its source. Write $\mu$ for the full $T$-character of $U|_S$. Every nonzero eigensection of $U$ has $\lambda$-weight at most $\langle\mu,\lambda\rangle$. If $U$ is nef, every fixed-component fiber weight has pairing at most that value and, for every $m\ge0$,
\begin{equation}\label{inj:extreme-index}
 [\chi_T(X,mU)]_{m\mu}=\chi(S,mU|_S).
\end{equation}
The trivial torus is allowed, with $S=X$ and $\mu=0$.
\end{lemma}
\begin{proof}
Bia\l ynicki-Birula's theorem \cite{BB1973} gives positive normal $\lambda$-weights at $S$. The first nonzero conormal jet of an eigensection subtracts a nonnegative sum of these weights from the line's fiber weight, proving the section bound. For an ample line, generation of a sufficiently high power gives the fixed-fiber weight bound. Apply this to $bU+A$, for an ample linearized $A$, and let $b\to\infty$ to obtain the bound for nef $U$.

To verify the full-character identity, evaluate the equivariant fixed-point formula \cite{AtiyahSegal1968,AtiyahSinger1968} at $a\lambda(t)$ as $t\to\infty$, divide by $t^{m\langle\mu,\lambda\rangle}$, and keep $a$ generic in $T$. The inverse normal Euler factors are
\[
 (1-a^{-w}t^{-\langle w,\lambda\rangle}e^{-x})^{-1}.
\]
They tend to one at $S$. Every other fixed component has a negative normal weight, hence an inverse factor tending to zero, and its line weight is no larger. Truncating in the ordinary cohomology degree gives the limit $a^{m\mu}\chi(S,mU|_S)$. Since the index is a finite Laurent polynomial, existence of this limit excludes higher $\lambda$-weights. The limit is precisely the sum of its terms on the weight hyperplane with pairing $m\langle\mu,\lambda\rangle$; varying $a$ identifies that sum with the single character $m\mu$ and the stated coefficient. For the trivial action the identity is immediate.
\end{proof}

\begin{theorem}[Exact attainment of source characters]
\label{inj:extremal-theorem}
Let $X$ be smooth connected projective with smoothly semipositive
$L=-K_X$ and with $H^0(X,\Omega_X^p)=0$ for all $p>0$.
Let an algebraic torus $T$ act with a linearized ample bundle.
For every generic integral cocharacter $\lambda$ with $X^\lambda=X^T$, let $S$ be its source and let $\mu$ be the full $T$-character of $L|_S$. Then for some $k_\mu>0$ there is a nonzero section of $k_\mu L$ of character $k_\mu\mu$. In particular a positive power of $L$ has a nonzero invariant section, for the natural linearization.
\end{theorem}

\begin{proof}
Take a generic integral cocharacter $\lambda$ with
$X^\lambda=X^T$.  A linearized projective embedding shows that
only finitely many walls must be avoided.  Bia\l ynicki-Birula's
attracting-cell theorem \cite{BB1973} gives a unique source
component $S$ with open attracting stratum as $t\to0$; all its
normal $\lambda$-weights are positive.  Projection of that stratum
to $S$ gives a dominant rational map $X\dashrightarrow S$.
Pullback of forms is injective and extends across the codimension-two
indeterminacy locus on smooth $X$.  Thus $S$ also has no
positive-degree holomorphic forms and
$\chi(S,\mathcal O_S)=1$.

Let $\mu$ be the full $T$-character of $L|_S$. Apply Lemma~\ref{lem:source-character} with $U=L$.

The right side of \eqref{inj:extreme-index} is polynomial in $m$
with value one at zero.  For infinitely many $m>0$, one fixed
cohomological degree has a nonzero weight-$m\mu$ class.
Equivariant hard Lefschetz with coefficient $(m+1)L$ produces,
for a fixed $p$, nonzero eigenforms
\[
 \alpha_m\in H^0(X,\Omega_X^p\otimes(m+1)L),
 \qquad\operatorname{wt}(\alpha_m)=m\mu.
\]
Use an invariant K\"ahler form and the maximal compact torus to
preserve all character spaces.

Use the determinant construction in Lemma~\ref{har:lem:forms}: the
generic span has rank $r$, and its saturated determinant is a linearized
line $M_0\subset\bigwedge^r\Omega_X^p$. Wedge forms with unbounded $m$
against $r-1$ fixed complementary basis forms. If the indices in one
such wedge are $m_1,\ldots,m_r$, its twist degree and character are
\[
 i=\sum_{j=1}^r(m_j+1),\qquad
 \sum_{j=1}^r m_j\mu=(i-r)\mu.
\]
Thus there are sections of $M_0+iL$ of weight $(i-r)\mu$ for unbounded
$i$. If $p>0$, the cotangent tensor theorem
\cite[Theorem 4.1]{LMPTX2023} makes $-M_0$ pseudoeffective;
for $p=0$, $M_0=\mathcal O_X$ as an ordinary line bundle.
Let $P$ be $L$ with its natural linearization shifted by $-\mu$,
and let $M$ be $M_0$ with its induced linearization shifted by
$r\mu$. These changes leave the ordinary line bundles unchanged,
so $-M$ is pseudoeffective in both cases above. The eigensections
just constructed are invariant sections of $M+iP$, because
\[
 (i-r)\mu+r\mu-i\mu=0.
\]
Theorem~\ref{inj:transfer} yields a nonzero section
\[
 v_\mu\in H^0(X,k_\mu L),\qquad
 \operatorname{wt}(v_\mu)=k_\mu\mu,
 \qquad k_\mu>0
\]
in the natural linearization.

There are finitely many source weights as the generic cocharacter
varies.  Zero lies in the convex hull of the actual section weights
$k_\mu\mu$.  Otherwise strict separation gives a generic integral
cocharacter pairing negatively with all of them.  At its own source,
the weight $\mu$ is a sum of tangent weights of positive pairing,
or the empty sum for a trivial action, contradicting strict
negativity.  The weights are integral, so a convex combination
giving zero may be chosen rational.  Clear its denominators and
multiply the corresponding powers of the nonzero $v_\mu$.
The product is a nonzero invariant section in positive degree.
For the trivial torus the same argument starts directly from
$\chi(X,\mathcal O_X)=1$ with $\mu=0$.
\end{proof}

\section{Singular injectivity and an Iitaka conversion}
\label{inj:injectivity}

The trivial-torus case of Theorem~\ref{inj:transfer} already converts
unbounded fixed-error sections into an anticanonical section. We give an
independent proof using singular injectivity and vanishing for all higher
direct images. On the Iitaka model $f:V\to Y$ constructed below, with
$V$ birational over $X$, the final interpolation retains
$f_*\mathcal O_V(K_{V/X})$ itself. The character-sensitive proof instead retains a corrected base
line with a fixed multiplier ideal. The resulting twisted-form conversion
will be applied after forms descend to a finite étale cover in
Section~\ref{har:subsec:formdescent}. The base metric need only have locally
bounded plurisubharmonic weights; the anticanonical metric upstairs
remains smooth and semipositive.

\begin{lemma}[Injectivity and base cohomology]\label{inj:leray}
Let $g:V\to Y$ be a morphism of smooth projective complex varieties, let $A$
be very ample on $Y$, and give $A$ a smooth metric with positive curvature
$\omega_A$.  Suppose a line bundle $N$ on $V$ has a singular Hermitian
metric $h$ such that
\[
 \mathcal I(h)=\mathcal O_V,
 \qquad i\Theta_h(N)\ge\epsilon g^*\omega_A
 \quad\text{for some }\epsilon>0.
\]
Then
\[
 H^i(Y,R^jg_*\mathcal O_V(K_V+N))=0
 \qquad(i>0,\ j\ge0).
\]
Disconnected smooth source varieties are allowed, with the hypotheses on
each component.
\end{lemma}

\begin{proof}
The input is the injectivity theorem of Fujino--Matsumura
\cite[Theorem A]{FM2016}.  On a compact K\"ahler manifold it applies to a
possibly singular coefficient metric and a smooth semipositive multiplying
line bundle $S$, provided the former curvature dominates a positive
multiple of the latter.  Multiplication by any nonzero section of $S$
is injective on canonical cohomology with the coefficient multiplier
ideal.  Here we use $S=g^*(mA)$; the comparison constant can be
$\epsilon/m$, and the ideal is trivial.

We include the reduction to base cohomology, following the slicing and
Leray argument in \cite[proof of Proposition 1.9, Step 3]{FM2016}.
Induct on $\dim Y$.  Choose $m$ large enough for Serre vanishing for all
the finitely many sheaves $R^jg_*\mathcal F$, where
$\mathcal F=\mathcal O_V(K_V+N)$.  A sufficiently general divisor
$C\in|mA|$ is smooth, has smooth inverse image $V_C$, contains no
associated point of these direct images, and contains the image of no
component of $V$.  Its restricted coefficient metric has trivial
multiplier ideal.  To see the last assertion without assuming analytic
singularities, pull the locally integrable metric coefficients to the
incidence variety of the basepoint-free system on $V$.  That incidence
variety is a projective-space bundle over $V$.  Fubini's theorem gives
integrability on almost every smooth member, simultaneously on a finite
coordinate cover.  This full-measure condition meets the required
algebraic open conditions.  Restricting the curvature inequality on
these members gives the same hypothesis over $C$.

Adjunction and multiplication by the cutting section give
\[
 0\longrightarrow\mathcal F\longrightarrow
 \mathcal F\otimes g^*\mathcal O_Y(mA)\longrightarrow
 \mathcal O_{V_C}(K_{V_C}+N|_{V_C})\longrightarrow0.
\]
The corresponding long direct-image sequence breaks into short exact
sequences: multiplication by the cutting equation is injective on every
$R^jg_*\mathcal F$, by the associated-point choice.  The induction
hypothesis on $C$ and Serre vanishing for the middle terms therefore
give $H^i(Y,R^jg_*\mathcal F)=0$ for $i\ge2$.

The Leray spectral sequence now has only columns zero and one.  For
each $j$ it identifies $H^1(Y,R^jg_*\mathcal F)$ with the first
filtration subspace of $H^{j+1}(V,\mathcal F)$.  After the twist by
$g^*(mA)$ this subspace maps to zero, since all positive base cohomology
of the twisted direct images vanishes.  Functoriality of the Leray
filtration and Fujino--Matsumura injectivity imply that the original
subspace is zero.  This proves the remaining degree $i=1$.
\end{proof}

\begin{proposition}[Conversion on an Iitaka base]\label{inj:ordinary-conversion}
Let $X$ be smooth connected projective over $\mathbb C$ and let
$L=-K_X$ carry a smooth semipositive Hermitian metric.  Suppose that
$F_0$ is a pseudoeffective line bundle and, for an unbounded set
$S\subset\mathbb Z_{>0}$, there are effective integral divisors
\[
 N_s\sim sL-F_0\qquad(s\in S).
\]
Then $H^0(X,mL)\ne0$ for some integer $m>0$.
\end{proposition}

\begin{proof}
Fix $a<b$ in $S$.  The affine relations
\begin{equation}\label{inj:affine-divisors}
 (s-a)N_b\sim(s-b)N_a+(b-a)N_s\qquad(s>b)
\end{equation}
will make the fiberwise section direction unique.  Put
$e=\kappa(X,N_b)$.  If $e=0$, the two effective divisors in
\eqref{inj:affine-divisors} are equal.  Letting $s$ tend to infinity
gives $N_b\ge N_a$, so $N_b-N_a\sim(b-a)L$ is effective.  If
$e=\dim X$, $bL=N_b+F_0$ is big, which also proves the assertion.
We treat $0<e<\dim X$.

Resolve a multiple system of $N_b$ attaining its Iitaka dimension and
take its Stein factorization.  After modifying the base, flattening
the main component, normalizing, and resolving, we have
\[
 V\longrightarrow W\longrightarrow X,
 \qquad f:V\to Y,
 \qquad f_W:W\to Y.
\]
Here $Y,V$ are smooth projective, $W$ is normal projective and birational
to $X$, $f_W$ is equidimensional, and $f$ has connected fibers.  The
normalization is finite and preserves the fiber-dimension bound from
flattening; the resolution $V\to W$ need not be equidimensional over
$Y$.  Equivalently the flat main component is obtained from the closure
of the family of generic fibers in the projective Hilbert scheme
\cite[Theorem~3.2]{Grothendieck1961Hilbert}.  Write
$\pi:V\to X$, $L'=\pi^*L$, and $E=K_{V/X}\ge0$.
For a very ample $A$ on $Y$ and some integer $t>0$,
\[
 t\pi^*N_b-f^*A\sim G\ge0.
\]
The initial moving polarization stays big after base modification, so
it dominates the chosen $A$ after increasing $t$.

For a geometric generic fiber $F$,
\begin{equation}\label{inj:rigidity}
 \kappa\bigl(F,(\pi^*N_b+E)|_F\bigr)=0.
\end{equation}
Indeed adding an effective exceptional divisor changes no global
sections of multiples of a pullback from the smooth $X$.  If two
sections on the generic fiber were independent, a sufficiently large
ample base twist would globalize two such sections.  Absorb that twist
using $t\pi^*N_b\ge f^*A$, up to linear equivalence.  Their nonconstant
fiberwise ratio, together with the ratios giving the base, would yield
image dimension at least $e+1$ for a multiple of $\pi^*N_b+E$.
The geometric generic fiber is integral because the Stein field is
relatively algebraically closed in characteristic zero.  Generic base
change justifies the globalization over the original function field.
This contradiction proves \eqref{inj:rigidity}.  The same assertion
holds without $E$.

Restriction of \eqref{inj:affine-divisors} to $F$ is therefore equality
of effective divisors.  Unboundedness of $S$ shows
$(\pi^*(N_b-N_a))|_F\ge0$.  If $d=b-a$, a rational section $z$ of
$dL$ with divisor $N_b-N_a$ is consequently regular on $F$.  For every
$k\ge0$,
\begin{equation}\label{inj:rank-one-all-k}
 h^0\bigl(F,(E+kdL')|_F\bigr)=1.
\end{equation}
The indicated effective section gives the lower bound.  For $k\ge1$,
multiply by the section of $k\pi^*N_a+(k-1)E$ to inject into the
sections of $k(\pi^*N_b+E)|_F$.  For $k=0$, multiply by the section
of $\pi^*N_b|_F$.  Equation \eqref{inj:rigidity} gives both upper
bounds.  This explicitly includes the zero exponent.

On the normal equidimensional model $W$, take for every prime $P\subset Y$
the minimum
\[
 c_P=\min_{Q\mapsto P}
       \frac{\operatorname{ord}_Q(z)}{
       \operatorname{ord}_Q(f_W^*P)}.
\]
Every vertical prime maps to a base prime and every denominator is
positive.  There are only finitely many nonzero $c_P$.  Subtracting
$f_W^*(\sum c_PP)$ leaves an effective rational Cartier divisor,
because the horizontal divisor of $z$ was already effective.
Clear denominators, replacing $d,z$ by a common power.  We obtain
\begin{equation}\label{inj:base-decomposition}
 dL'\sim H+f^*B,
 \qquad H\ge0,
\end{equation}
with $H$ pulled back from an effective Cartier divisor on $W$ and
$B$ integral on $Y$.  Above each base prime at least one component
has coefficient zero in $H$.

Proposition~\ref{har:prop:maximum} now applies with the ordinary line $P=L$, the rational section $z$, and the normal equidimensional model $W$. Equation~\eqref{inj:rank-one-all-k} supplies rank one for every $k\ge0$; the vertical minima giving \eqref{inj:base-decomposition} agree with the proposition's normalization. Hence $B$ has a continuous semipositive metric. In a local frame $b$, with $\sigma=s_Hf^*b$ and $h_0=\pi^*h_L$, its weight is
\begin{equation}\label{inj:max-weight}
 \phi_B=-\log\max_{V_y}|\sigma|_{h_0^d}^2.
\end{equation}
Only local boundedness of this metric is needed in the following injectivity argument.

It remains to use the zero exponent rather than merely asymptotic
positivity.  Since $F_0$ is pseudoeffective, the ample domination above
gives $nL'-f^*A$ pseudoeffective for some $n>0$.  Its positive current
yields a singular metric on $L'$ with curvature dominating a positive
multiple of $f^*\omega_A$.  Mix its weights with the smooth semipositive
weights using a sufficiently small positive coefficient.  Skoda's
small-exponent integrability \cite[Lemma~5.6(a)]{Demailly2011}, in
the squared-norm convention $e^{-\phi}$, and a finite covering of $V$ make the
new multiplier ideal trivial while retaining a strictly positive base
curvature coefficient.  No analytic-singularity assumption is used.
For $k\ge0$ put $C_k=L'+kf^*B$ and add $kf^*\phi_B$ to the coefficient
weights. Local boundedness keeps the multiplier ideal trivial and
preserves the positive base-curvature lower bound. Since
$K_V+C_k=E+kf^*B$, Lemma~\ref{inj:leray} and the projection formula give
\[
 H^i\bigl(Y,f_*\mathcal O_V(E)\otimes\mathcal O_Y(kB)\bigr)=0
 \quad(i>0,\ k\ge0).
\]
Riemann--Roch makes the Euler characteristic a polynomial in $k$.
At $k=0$ it equals $h^0(V,E)=1$.  It is therefore nonzero at some
positive integer $k$.  A section of $E+kf^*B$, multiplied by $s_H^k$,
gives a section of $E+kdL'$.  Since
$\pi_*\mathcal O_V(E)=\mathcal O_X$, it yields a nonzero section of
$kdL$ on $X$.
\end{proof}

\begin{corollary}[Twisted-form input]\label{inj:forms-input}
Under the smooth projective and smooth anticanonical semipositivity
hypotheses of Proposition~\ref{inj:ordinary-conversion}, suppose
$H^0(X,\Omega_X^p\otimes mL)\ne0$ for an unbounded set of positive
integers $m$ and a fixed $p\ge0$.  Then a positive power of $L$ has
a nonzero section.
\end{corollary}
\begin{proof}
For $p=0$ the assertion is immediate.  Otherwise take the stable generic
span of the directions of the twisted forms in $\Omega_X^p$, of rank
$r>0$, and saturate its determinant in
$\bigwedge^r\Omega_X^p\subset(\Omega_X^1)^{\otimes pr}$.
A saturated rank-one subsheaf of a locally free sheaf on smooth $X$
is reflexive and hence invertible; call it $M$.  Wedge an arbitrarily
far-tail member with suitable $r-1$ members of a fixed generic basis.
This gives nonzero sections of $M+sL$ for unbounded sums $s$.
The cotangent-tensor generic-nefness theorem
\cite[Theorem 4.1]{LMPTX2023} applies because $-K_X$ is nef and
proves that $-M$ is pseudoeffective.  Apply the proposition with
$F_0=-M$ and the divisors of these sections.
\end{proof}

For the direct twisted-form descent in
Section~\ref{har:subsec:formdescent}, the forms in this
corollary must first descend to a smooth projective finite \'etale
cover of the original variety.  The conversion asserts ordinary
sections on that cover.  A norm then descends a section of a further
positive power.  It asserts neither equivariance of the conversion nor
a uniform bound on the positive exponent.

\section{Harmonic restriction with a continuous base metric}\label{har:sec:harmonic}

The coefficient on the total space in the decomposition theorem below
always has a smooth metric. Continuity enters only through a line bundle
pulled back from the base. We give the approximation and restriction
arguments because the decomposition theorem by itself is not a vanishing
theorem.

\begin{lemma}[Uniform smoothing]\label{har:lem:smoothing}
Let \(Q\) be a line bundle on a smooth projective \(Y\), with a positive
continuous Hermitian metric whose local weights are plurisubharmonic.
For a fixed Kähler form \(\eta\), there are smooth metrics whose weights
converge uniformly to these weights and whose curvatures are at least
\(-\varepsilon_j\eta\), with \(\varepsilon_j\to0\).
\end{lemma}

\begin{proof}
This is the regularized-maximum construction of Richberg
\cite{Richberg1968}; see also \cite[Theorem~I.5.21]{Demailly2012}. The small negative curvature allowance permits its
use for semipositive, rather than strictly positive, weights. Write a
weight as a smooth reference weight plus a global continuous function
\(u\), with background curvature \(\gamma\). For fixed \(\varepsilon>0\),
use background \(\gamma+\varepsilon\eta\), so that the original potential
has a strict local margin. On finitely many enlarged coordinate patches,
add a smooth local potential of this background, convolve, and subtract
that potential. Take inner patches which still cover \(Y\). Add smooth
bumps zero on the inner patches and a small negative constant near the
outer boundaries, choosing their second derivatives smaller than the
strict margin and the convolution errors much smaller than the bump
heights. A regularized maximum of the active candidates is smooth: a
candidate near its outer boundary is uniformly below a candidate from
an inner patch and does not affect the maximum. Convexity, monotonicity,
and translation equivariance of regularized maximum preserve the
background lower bound and the transition rule for weights. Let the
uniform error and then \(\varepsilon\) tend to zero. Equivalently, one
may smooth the strictly positive metric on \(NQ+H\), with \(H\) positive,
subtract the weight of \(H\), and divide by \(N\).
\end{proof}

\begin{theorem}[Kernel-condition vanishing]\label{har:thm:kernel}
Let $f:V\to Y$ be a surjective morphism between smooth connected projective complex varieties. Let $D_V$ be a line bundle with a smooth semipositive Hermitian metric, whose curvature is $\theta$. Let $Q$ be a line bundle on $Y$ with a positive continuous Hermitian metric having plurisubharmonic weights.
Assume that on a nonempty ordinary open subset of the submersion locus,
\begin{equation}\label{har:eq:kernel}
 \ker\theta\subseteq\ker df.
\end{equation}
Then, for every \(k\ge0\), \(q\ge0\), and \(p>0\),
\begin{equation}\label{har:eq:vanishing}
 H^p\bigl(Y,R^q f_*(K_V+D_V)\otimes Q^k\bigr)=0.
\end{equation}
In applications below \(D_V\) is the pullback of a smoothly semipositive
anticanonical line, and the particular diagram supplies
\eqref{har:eq:kernel}; the rank-open condition is not inferred from nefness.
\end{theorem}

\begin{proof}
Fix \(k\). We will show that every nonzero cohomology class on the
source restricts nontrivially to some general fiber. Smooth-coefficient
decomposition will then give the asserted vanishing. Choose a Kähler
form \(\omega\) on \(V\) and a Dolbeault class
\(c\in H^\ell(V,K_V+A)\), where \(A=D_V+kf^*Q\) and \(n=\dim V\).
Lemma~\ref{har:lem:smoothing} gives product metrics \(h_j\to h\) uniformly
on \(A\), with curvature at least \(\theta-o(1)\omega\). Let \(u_j\)
be the harmonic \((n,\ell)\)-representative of \(c\) for \(h_j\).
Harmonic norm minimization, comparison with one fixed smooth closed
representative, and uniform metric equivalence give a fixed \(L^2\) bound.
The Bochner--Kodaira identity gives
\begin{equation}\label{har:eq:energy}
 \|(D'_j)^*u_j\|_{h_j}^2
 +\langle[\theta,\Lambda]u_j,u_j\rangle_{h_j}
 \le o(1)\|u_j\|_{h_j}^2.
\end{equation}
Here curvature commutators in bidegree \((n,\ell)\) are monotone in the
curvature form: their eigenvalues are sums of the curvature eigenvalues
on the antiholomorphic index set.

Let \(v_j=*u_j\), absorbing a fixed nonzero dimensional constant so that
\(u_j=\omega^\ell\wedge v_j\). The complex-linear star acts on form
indices, carrying the coefficient line along. Metric compatibility gives
\((D'_j)^*=\pm*\bar\partial*\) and
\(\bar\partial_j^*=\pm*D'_j*\). Thus
\[
 \bar\partial v_j\longrightarrow0\text{ in }L^2,
 \qquad D'_jv_j=0.
\]
These are statements about a fixed holomorphic bundle; the first
operator does not differentiate the varying metric. Passing to a weak
\(L^2\) subsequence gives a distributionally holomorphic, hence smooth
holomorphic, \(v\). Since \(\omega\) is Kähler, the form
\(u=\omega^\ell\wedge v\) is \(\bar\partial\)-closed and represents
\(c\). Indeed projection to harmonic forms for one fixed smooth
reference metric is a bounded finite-rank operator, has the same value
on every \(u_j\), and commutes with this weak limit. The nonnegative
square root of \([\theta,\Lambda]\) annihilates the limit by \eqref{har:eq:energy}.
Pointwise diagonalization therefore shows that every positive-curvature
holomorphic covector is a wedge factor of \(v\).

The approximation also gives a Chern equation before any use of the
kernel condition. In any local holomorphic frame of \(A\), let
\(h_j\) denote its squared norm. The equation \(D'_jv_j=0\) reads
\(\partial(h_jv_j)=0\). Uniform convergence of \(h_j\) and weak
\(L^2\) convergence suffice to pass it to distributions:
\begin{equation}\label{har:eq:distribution}
 \partial(hv)=0.
\end{equation}
No derivative convergence of metric weights is used.

To detect the class on a fiber, we next show that \(v\) contains all
base canonical covectors as wedge factors. On the open set in
\eqref{har:eq:kernel}, every pulled-back base covector annihilates
\(\ker\theta\), so its wedge with \(v\) vanishes. The holomorphic
bundle map
\[
 f^*\Omega_Y^1\longrightarrow\Omega_V^{n-\ell+1}\otimes A,
 \qquad \alpha\longmapsto (df)^*\alpha\wedge v,
\]
therefore vanishes everywhere. Over the smooth fibration locus, in
base coordinates \(y_1,\ldots,y_b\), this gives
\begin{equation}\label{har:eq:basefactor}
 v=f^*(dy_1\wedge\cdots\wedge dy_b)\wedge w.
\end{equation}
The relative restriction \(w_y\) is unambiguously defined after choosing
the base canonical frame. If \(n-\ell<b\), then \(v=0\) already.

The limiting metric is the smooth metric of \(D_V\) times a positive
continuous function pulled back from the base. Because of
\eqref{har:eq:basefactor}, only vertical differentiations survive in
\eqref{har:eq:distribution}.
The continuous base function is constant under these differentiations;
the remaining expression is smooth in the vertical variables and
continuous altogether. It vanishes pointwise and gives
\(D'_{D_V|V_y}w_y=0\) on every smooth fiber under consideration.

The restriction in cohomology is restriction of a \((0,\ell)\)-form
\emph{valued in} \(K_V+A\). Use
\(K_V|_{V_y}=K_{V_y}\otimes K_Y|_y\) to remove the constant canonical
base factor. It is not the ordinary pullback of an \((n,\ell)\)-form to
a lower-dimensional fiber. The restricted representative is a nonzero
constant times
\[
 \omega_y^\ell\wedge w_y,
 \qquad \omega_y=\omega|_{V_y}.
\]
It is harmonic on the compact smooth fiber: it is \(\bar\partial\)-closed,
and its star has zero Chern derivative. Thus a nonzero \(v\) gives a
nonzero restriction class on some general fiber. The Leray edge map
\begin{equation}\label{har:eq:edge}
 H^\ell(V,K_V+D_V+kf^*Q)
 \longrightarrow H^0(Y,R^\ell f_*(K_V+D_V)\otimes Q^k)
\end{equation}
is injective, since zero edge image implies zero restriction on the
general smooth base-change locus.

Finally apply Fujisawa's derived decomposition
\cite[Theorem~1]{Fujisawa2015}, obtained from Takegoshi's harmonic
representatives, to the \emph{smooth} coefficient \(D_V\):
\[
 Rf_*(K_V+D_V)\simeq\bigoplus_q R^q f_*(K_V+D_V)[-q].
\]
Tensor the resulting derived isomorphism by the line \(Q^k\) on \(Y\)
and use the projection formula. It gives
\[
 h^\ell(V,K_V+D_V+kf^*Q)
 =\sum_{p+q=\ell}h^p(Y,R^qf_*(K_V+D_V)\otimes Q^k).
\]
Injection \eqref{har:eq:edge} bounds the left side by the summand \(p=0\).
All remaining summands vanish. For total degrees exceeding \(n\), the
left side is zero directly. This proves every case of
\eqref{har:eq:vanishing}, including the twist \(k=0\).
\end{proof}

\begin{remark}[The vertical-polyvector proof]\label{har:rem:polyvector}
There is a second way to obtain the restriction step, useful when only
uniform comparability, rather than convergence of metrics, is known.
Let \(\rho\) be a smooth semipositive real \((1,1)\)-form on \(V\). Suppose a line \(A\) has uniformly comparable smooth metrics of curvature
\(\ge c\rho-\varepsilon_j\omega\), where \(c>0\) is fixed and \(\varepsilon_j\to0\), on a smooth projective \(V\).
The weak-limit argument through \eqref{har:eq:energy} gives a holomorphic
\(v\in H^0(\Omega_V^{n-\ell}\otimes A)\) representing each adjoint class
by \(\omega^\ell\wedge v\). Under
\(\Omega_V^{n-\ell}\otimes A=\bigwedge^\ell T_V\otimes(K_V+A)\),
it corresponds to a holomorphic polyvector whose indices lie in
\(\bigwedge^\ell\ker\rho\). If \(\ker\rho\subset\ker df\) on a
nonempty open, this polyvector is vertical throughout the submersion
locus by holomorphy. Lowering its vertical indices and restricting uses
exactly the induced fiber Kähler metric, with the canonical base factor
held constant. Suppose on each general fiber there is a smooth metric
on \(A|_{V_y}\) with curvature a positive multiple of \(\rho|_{V_y}\).
The corresponding holomorphic \((\dim V_y-\ell,0)\)-form \(v_y\)
contains all the positive-curvature covectors. Hence its curvature
commutator is zero. Since a holomorphic form of antiholomorphic degree
zero is \(\bar\partial\)-harmonic, Bochner--Kodaira gives \(D'v_y=0\).
Its Lefschetz image is fiber-harmonic. This proves the same detection
and edge injection without passing a Chern equation through the limit.
The smooth-coefficient splitting is still a separate requirement.
\end{remark}

\begin{proposition}[Pseudoeffective domination variant]\label{har:prop:dominationvanish}
In Theorem~\ref{har:thm:kernel}, replace the kernel hypothesis by the
assumption that \(bD_V-f^*H\) is pseudoeffective for some \(b>0\) and
some ample line \(H\) on \(Y\). Then the same vanishing holds.
\end{proposition}

\begin{proof}
Keep the uniformly convergent metrics and the notation of the proof of
Theorem~\ref{har:thm:kernel}. For \(\ell>0\), put
\[
 S=i^{(n-\ell)^2}h\,v\wedge\overline v,
 \qquad T=S\wedge\omega^{\ell-1},
\]
with coefficient pairing understood. The distributional Chern equation
and holomorphy make \(T\) a continuous closed positive
\((n-1,n-1)\)-form. The curvature-energy equality gives
\(\theta\wedge T=0\). Choose a positive curvature form \(\eta\) of
\(H\). A positive current in the class of
\(bD_V-f^*H\) pairs nonnegatively with \(T\), since positive currents
have order zero. To compute this pairing cohomologically, apply de Rham
regularization, which commutes with \(d\), to approximate \(T\)
uniformly by smooth closed forms. For each approximant the pairing
agrees with that against the cohomologous smooth form
\(b\theta-f^*\eta\); uniform convergence passes this equality to
\(T\). It follows that
\[
 0\le\int_V(b\theta-f^*\eta)\wedge T
       =-\int_V f^*\eta\wedge T\le0.
\]
Pointwise positivity forces \(f^*\eta\wedge T=0\); diagonalization gives
the full base wedge factor in \(v\). The distributional fiber equation,
canonical-valued restriction, edge injection, and smooth-coefficient
splitting proceed exactly as in that proof. Only positive total degrees
are needed to kill a term of positive base degree.
\end{proof}

\begin{proposition}[Euler interpolation on a birational model]\label{har:prop:euler}
Let \(\pi:V\to X\) be birational between smooth connected projective
varieties, let \(L=-K_X\) be smoothly semipositive, put
\(D_V=\pi^*L\), and write \(E=K_{V/X}=K_V+D_V\).
Let \(f:V\to Y\) be surjective, with \(Y\) smooth connected
projective. Suppose a continuously semipositive line \(Q\) on \(Y\)
and an injection \(f^*Q\to bD_V\), \(b>0\), satisfy either vanishing
criterion above. Then \(H^0(X,kbL)\ne0\) for some \(k>0\).
If a torus acts, trivially on \(Y\) and \(Q\), and all maps and the
injection are equivariant, this section can be taken invariant.
\end{proposition}

\begin{proof}
Let \(\mathcal G=f_*\mathcal O_V(E)\), or its invariant summand in the
equivariant case. A torus over a trivially acted-on base has coherent
weight summands, so this is a direct summand. Vanishing gives
\(h^0(Y,\mathcal G\otimes Q^k)=\chi(Y,\mathcal G\otimes Q^k)\) for every
\(k\ge0\). The right side is polynomial in \(k\) by Riemann--Roch for
coherent sheaves. At zero it is one: \(\pi_*\mathcal O_V(E)=\mathcal O_X\)
and the exceptional canonical section is invariant. The polynomial is
therefore nonzero at some positive integer. Multiply the resulting
section of \(E+kf^*Q\) by the \(k\)-th power of the injection and push
through \(\pi\). The effective exceptional divisor \(E\) changes neither
sections nor their characters, by normality of \(X\).
\end{proof}

\section{Invariant Euler nonvanishing by a fixed invariant system}\label{har:sec:eulerroute}

We now give a harmonic proof of Corollary~\ref{thm:main}. Starting with
the determinant sections, we choose the base from powers of one fixed
invariant line. Maximality of its image will force the invariant
adjoint spaces on the generic fiber to have rank one. The determinant
intersection identity will supply the curvature-kernel condition.
Together these properties give the continuously semipositive base line
and the vanishing at twist zero required by
Proposition~\ref{har:prop:euler}. The premise here is the nonzero
invariant Euler characteristic; vanishing of every higher
structure-sheaf cohomology group is not required.

\begin{proof}[Harmonic proof of Corollary~\ref{thm:main}]
Use Lemma~\ref{har:lem:forms}, retaining the cases which did not already
produce a section, and write \(N_i=M+a_iL\), \(s_i\) for its invariant
sections. Write \(n=\dim X\), let \(\theta\) be the averaged smooth
semipositive curvature, let \(r\) be its maximal rank, and fix a
K\"ahler form \(\omega\). Choose indices \(a_-<a\) and keep unbounded indices above
\(a\). Put \(N=M+aL\). Among rational maps given by all invariant
sections of positive powers of \(N\), choose one with maximal image
dimension. Resolve its base ideal equivariantly and take Stein
factorization. We obtain a smooth projective \(X_0\), a birational
\(\pi_0:X_0\to X\), and \(f_0:X_0\to Y_0\) with connected fibers and
normal projective base, with
\begin{equation}\label{har:eq:iitaka}
 d\pi_0^*N=f_0^*H_0+G_0,\qquad G_0\ge0,
\end{equation}
where \(H_0\) is ample. The torus acts trivially on \(Y_0,H_0\), and
the defining section of \(G_0\) is invariant.

The invariant sections of \(j\pi_0^*N\) on the generic fiber form a
one-dimensional vector space over \(\mathbb C(Y_0)\) for every \(j>0\).
Indeed finitely many generic-fiber sections extend rationally with
vertical poles. A section of a sufficiently high power of \(H_0\)
clears those poles; multiplication by the corresponding power of the
section of \(G_0\) gives invariant sections of a higher power of
\(\pi_0^*N\), and hence of \(N\) on \(X\). If two original sections
were independent, their ratio would not lie in \(\mathbb C(Y_0)\).
This field is relatively algebraically closed by Stein factorization,
so that ratio would be transcendental over it. Combine it with the
original image coordinates, by products to put all ratios in one
power. The image dimension increases, contradicting maximality.

For each \(a_i>a\) there are positive integers \(p,q\) such that
\(pa_i+qa_-=(p+q)a\). The invariant section
\(s_i^p s_-^q\) is thus generically proportional to \(s_a^{p+q}\).
Its effective generic-fiber divisor shows that all horizontal
components of \(\operatorname{div}(\pi_0^*s_i)\) lie in the fixed finite
support of \(\operatorname{div}(\pi_0^*s_a)\). Extract a subsequence of
their integer multiplicity vectors that is componentwise nondecreasing.
For two indices \(a_j>a_i\) in it, put
\[
 b=a_j-a_i>0,\qquad s=\pi_0^*s_j/\pi_0^*s_i.
\]
This is an invariant rational section of \(b\pi_0^*L\), regular on the
generic fiber. The relative canonical divisor
\(E_0=K_{X_0/X}\) has its horizontal support in the same fixed support:
\(\pi_0\) can be chosen to be an isomorphism outside the invariant base
ideal, which lies in the zero locus of \(s_a\). Fix \(k\ge0\), and work on the generic fiber.  The effective
divisors of \(s_{E_0}\) and \(s\) have support in the divisor of
\(\pi_0^*s_a\).  All coefficients are finite, so a sufficiently large
integer \(h\) satisfies
\[
 h\operatorname{div}(\pi_0^*s_a)
 \ge \operatorname{div}(s_{E_0})+k\operatorname{div}(s).
\]
The quotient \((\pi_0^*s_a)^h/(s_{E_0}s^k)\) is therefore a
regular invariant section of the difference line.  Multiplication
by it injects the entire invariant adjoint section space into the
one-dimensional invariant space for \(h\pi_0^*N\).  Consequently
\begin{equation}\label{har:eq:invrank}
 \dim_{\mathbb C(Y_0)}H^0((X_0)_\eta,E_0+kb\pi_0^*L)^T=1
 \qquad(k\ge0).
\end{equation}
If \(\dim Y_0=0\), the section \(s\) is already regular on \(X_0\)
and descends to the required section on \(X\).

Otherwise apply Lemma~\ref{val:flat-model} to this fibration and resolve
its normal equidimensional model:
\[
 V\longrightarrow W\longrightarrow X_0,\qquad g:W\to Y,
 \qquad f:V\to Y.
\]
Write \(\mu:W\to X\) and \(\pi:V\to X\) for the composed birational maps. We also use \(f_0\) for its pullback to \(V\). Here \(Y\) is smooth projective, \(W\) is normal equidimensional over
\(Y\), and \(V\) is smooth. All constructions are equivariant, with
trivial base action. The ranks
\eqref{har:eq:invrank} persist: a further effective birational canonical
divisor has trivial pushforward also on the generic smooth fiber.

It remains to relate this algebraic fibration to the curvature. The
kernel condition for \(D_V=\pi^*L\) follows from
Lemma~\ref{har:lem:null-divisors} and \eqref{har:eq:iitaka}. If \(r<n\),
\[
 \int_V (\pi^*\theta)^r\wedge f_0^*\eta_{H_0}
                       \wedge(\pi^*\omega)^{n-r-1}=0.
\]
Indeed the whole class \(dN\) has zero intersection, and both the
moving and effective fixed parts have nonnegative intersection. On
an ordinary open where \(\pi\) is an isomorphism, \(\theta\) has rank
\(r\), and the base-to-image map has injective differential, positivity
forces the horizontal form to annihilate \(\ker\pi^*\theta\). Its
kernel there is \(\ker df\). Thus
\(\ker\pi^*\theta\subset\ker df\). These algebraic generic conditions
meet the nonempty ordinary maximal-rank open. For full curvature rank
the kernel inclusion holds directly.

Proposition~\ref{har:prop:maximum} normalizes \(s\) to
\(g^*Q\to b\mu^*L\), with effective zero divisor \(G\), and gives a
continuous semipositive metric on \(Q\). Theorem~\ref{har:thm:kernel}
and Proposition~\ref{har:prop:euler} now produce an invariant section
of a positive power of \(L\).

\end{proof}

\subsection{An alternative proof that the maximum stays positive}

The preceding proof is complete. We record a second argument for the
positivity of the fiberwise maximum on its model. Instead of extending
the maximum weight across a codimension-two set, we first prove
nefness by stabilizing the sheaves of allowed poles. Blowing up a
hypothetical zero fiber then contradicts nefness. This argument uses
harmonic vanishing and is independent of the maximum-principle step
in Proposition~\ref{har:prop:maximum}.

\begin{proposition}[Pole-sheaf stabilization on the invariant Iitaka model]\label{har:prop:pole-stabilization}
Let $V\to W\to X$, $g:W\to Y$, $f:V\to Y$, and $D_V=\pi^*(-K_X)$ be the normal equidimensional model and smooth resolution in the preceding invariant Iitaka construction. Put $E=K_{V/X}$. Retain its invariant adjoint rank-one condition \eqref{har:eq:invrank} and its curvature-kernel inclusion, and write $bD_V=G+f^*Q$ after vertical normalization, with $G$ pulled back from the normal model and omitting a divisorial lift of each base prime. Then the increasing invariant pole sheaves
$\mathcal P_k=(f_*\mathcal O_V(E+kG))^T$, viewed as rational functions through $s_Es_G^k$, stabilize. The line $Q$ is nef and $G$ contains no entire fiber. These conclusions give a second proof that the fiberwise maximum defines a positive continuous metric.
\end{proposition}
\begin{proof}
On \(V\), form
\[
 \mathcal P_k=(f_*\mathcal O_V(E+kG))^T\qquad(k\ge0),
\]
viewed as sheaves of rational functions on \(Y\) by division by
\(s_Es_G^k\). They have rank one, contain \(\mathcal O_Y\), and increase
with \(k\). For each prime \(P\) of \(Y\), take a lift \(\Gamma\)
omitted by \(G\). Regularity of \(f^*a\,s_Es_G^k\) gives
\[
 \operatorname{ord}_{\Gamma}(f^*P)\operatorname{ord}_P(a)
      +\operatorname{ord}_{\Gamma}E\ge0.
\]
The possible poles are therefore bounded independently of \(k\) by
one divisor \(C\). Normality embeds all \(\mathcal P_k\) into
\(\mathcal O_Y(C)\), so Noetherianity gives a stable value
\(\mathcal P_\infty\).

Take a very ample \(H\) on \(Y\). Apply
Theorem~\ref{har:thm:kernel}, with trivial base twist, to the smooth coefficient
\((1+kb)D_V+lf^*H\), \(k,l\ge0\). Its curvature dominates the original
curvature, so its kernel is still vertical on the same open set.
Since its adjoint line is \(E+kG+f^*(kQ+lH)\), projection to
invariants and the projection formula give
\[
 H^i(Y,\mathcal P_k\otimes(kQ+lH))=0\qquad(i>0).
\]
For large \(k\), Castelnuovo--Mumford regularity makes
\(\mathcal P_\infty\otimes(kQ+(\dim Y)H)\) globally generated. Pull
back to the normalization of any curve and remove torsion. The fixed
sheaf has positive rank, so its degree plus \(k\) times that rank times
\(\deg Q\) is nonnegative for unbounded \(k\). Thus \(Q\) is nef.

If \(G\) contained the entire fiber over \(y\in Y\), blow up \(y\)
(the curve case is already excluded by normalization). The normalized
main transform remains equidimensional because it lies in a base
change with the same fiber-dimension upper bound. Recompute the
vertical minima on this model. The new rational base divisor is the
pullback of \(Q\) plus a strictly positive multiple of the exceptional
divisor: every divisorial lift of its generic point lies over the
old fiber and therefore has positive order in the pulled-back \(G\).
The rank and kernel arguments just given still apply, so the new
divisor is nef after clearing denominators. Its degree on a line in
the exceptional projective space is negative, a contradiction. Hence
no entire fiber lies in \(G\). Properness and openness give a positive
continuous maximum, and the moment calculation
\eqref{har:eq:moments} proves semipositivity without the codimension-two
extension argument. This gives the same conclusion by a distinct
pole-sheaf stabilization proof.
\end{proof}

\section{Two invariant fields giving the harmonic diagram}\label{har:sec:fields}

The preceding proof chooses a base from powers of one fixed line.
Here we give two alternatives. In the first, finite rank of the
linearized Picard group controls the horizontal divisors of the
determinant sections. In the second, equal-degree products force
their horizontal orders to be affine in the exponent. Both produce
a relative anticanonical section, rank-one invariant adjoint spaces,
and a vertical curvature kernel, so that
Propositions~\ref{har:prop:maximum} and~\ref{har:prop:euler} apply.

Throughout this section assume \(H^i(X,\mathcal O_X)=0\) for \(i>0\),
and let \(T\) act algebraically on the smooth connected projective
\(X\). Put \(n=\dim X\) and \(L=-K_X\), choose a compact-torus-invariant smooth
semipositive metric, and write \(\theta\) for its curvature and
\(r\) for its maximal rank. Apply Lemma~\ref{har:lem:forms}. If it
already produces an invariant anticanonical section, there is nothing
to construct; this includes the full-rank case. Otherwise write its
determinant sections as \(s_i\in H^0(X,N_i)^T\), where
\(N_i=M+a_iL\) and \(a_i\to\infty\).
Lemma~\ref{har:lem:null-divisors} gives their zero intersection with
\([\theta]^r[\omega]^{n-r-1}\). If \(r=0\), their effective zero
divisors have zero K\"ahler degree and hence vanish; a ratio of two
sections gives the conclusion directly. We therefore work with
\(0<r<n\).

\subsection{Invariant functions with curvature-null poles}\label{har:subsec:small}

Call a prime divisor small if \(\theta^r\) vanishes on its smooth
locus. An effective divisor has small support exactly when its
intersection with \([\theta]^r[\omega]^{n-r-1}\) vanishes.
Smallness persists under pullback to smooth birational
models and for exceptional divisors centered in small supports. To see
this even over a singular locus, resolve the original prime, take a
dominating component of the fiber product with the proposed source,
and test the pulled-back form on a smooth dominating space; continuity
extends its zero identity.

Let \(\mathcal K\) consist of zero and the invariant rational functions
with small pole support. It is a field: sums and products introduce no
new pole support, while zeros and poles of a rational function are
linearly equivalent effective divisors, so the zero support is small
whenever the pole support is. This also handles inversion. It is a
finitely generated intermediate field in \(\mathbb C(X)\). Explicitly,
take a transcendence basis of \(\mathcal K\), extend it to one of
\(\mathbb C(X)\), and note that adjoining the complementary independent
variables preserves the degrees of finite algebraic subextensions.
The finite degree of the full function field over the larger purely
transcendental field bounds those degrees, making the algebraic part
finite.

Use finitely many generators to map to a product of projective lines,
take the image \(Y_0\), and resolve equivariantly:
\(\pi_0:X_0\to X\), \(f_0:X_0\to Y_0\). The base action is trivial.
Indeterminacy lies in the small pole supports, so the resolution may
be chosen to be an isomorphism off them. The exceptional canonical
divisor \(E_0\) is small. A pullback of an ample line \(H_0\) on
\(Y_0\) has a small effective representative, built from the pole
divisors of the generators and exceptional corrections.

The choice of field now controls horizontal divisors. There are only
finitely many invariant small primes on \(X_0\) horizontal over
\(Y_0\). Indeed \(\operatorname{Pic}^0(X_0)=0\) by birational
invariance, and the kernel of the linearized Picard group over the
ordinary Picard group consists of characters. Thus the linearized
Picard group has finite rank. Infinitely many such primes would give,
after clearing torsion, a nonzero integral relation among their
natural linearized divisor lines. The corresponding quotient of
canonical divisor sections is an invariant rational function whose
divisor is that nonzero horizontal combination. Its poles downstairs
are small, so it belongs to \(\mathcal K=\mathbb C(Y_0)\).
But a function pulled back from \(Y_0\) has vertical divisor, a
contradiction.

The determinant divisors chosen above are small by
Lemma~\ref{har:lem:null-divisors}, so their horizontal supports belong
to this finite set. Order
the multiplicities along an infinite subsequence and take a ratio.
We obtain an invariant rational section \(s\) of \(b\pi_0^*L\),
\(b>0\), regular on the generic fiber. For every \(k\ge0\), both
\[
 (f_{0*}\mathcal O_{X_0}(kb\pi_0^*L))^T,
 \qquad (f_{0*}\mathcal O_{X_0}(E_0+kb\pi_0^*L))^T
\]
have rank one. They contain \(s^k\) and \(s_{E_0}s^k\). Any other
invariant generic section has its vertical poles cleared by a power
of \(H_0\). Its effective zero divisor is small, since its class has
zero intersection with the semipositive test forms, using smallness of
\(E_0,H_0\) and \(\theta^{r+1}=0\).
The ratio of two such sections therefore has small poles and is in
\(\mathcal K\). They are dependent over the base field, as required.

The kernel condition has a direct differential proof. Every level
divisor of a generating function is small. At a regular level point
in the rank-\(r\) open, its tangent hyperplane must contain
\(\ker\theta\): a hyperplane not containing that kernel would surject
onto the rank-\(r\) quotient, so \(\theta^r\) would not vanish on it.
Thus \(\ker\theta\subset\ker df_0\) on a nonempty ordinary open.
After an equidimensional normal model and a smooth resolution, all
these properties persist. Proposition~\ref{har:prop:maximum} gives the
continuous base metric and Proposition~\ref{har:prop:euler} gives an
invariant anticanonical section. No separate relative-algebraic-closure
assumption on the small-pole field is needed. The invariant adjoint
rank at \(k=0\) already forces geometric connectedness of the generic
fiber, as explained in Proposition~\ref{har:prop:maximum}.

\subsection{Ratios from a semigroup of determinant lines}\label{har:subsec:semigroup}

Retain the hypotheses and determinant sections fixed at the start of
the section. This time we impose all equal-degree ratios at once;
their relations will give a relative section without first listing
the horizontal primes. Form the semigroup of linearized line bundles generated by the
\(N_i\), and the field generated by ratios of invariant sections of
the same semigroup element. Choose a transcendence basis among these
ratios. Products put their numerators and denominators into one
invariant linear system of a semigroup element \(N_0\). Normalize its
image in the relative algebraic closure of its function field in
\(\mathbb C(X)\), and resolve the base to obtain \(Y\). Every original
ratio is algebraic over the chosen transcendence basis, so belongs
to this finite relative algebraic closure. The connected torus
acts trivially on it: it fixes the smaller field and has no nontrivial
finite quotient. Thus the base action is trivial.

After the Hilbert-family modification, take the normal equidimensional
main graph \(W\to Y\) and its smooth resolution \(V\). Write
\(\pi:V\to X\), \(f:V\to Y\), \(D_V=\pi^*L\), \(E=K_{V/X}\).
The pullback \(A\) of the image hyperplane bundle to \(Y\) is big,
and the common factor of the pulled-back coordinates is regular,
giving an invariant effective difference
\(\pi^*N_0-f^*A\). Hence for an ample \(H\) on \(Y\),
\begin{equation}\label{har:eq:absorb}
 a\pi^*N_0-f^*H\quad\hbox{is invariantly effective for some }a>0.
\end{equation}
The intersection identity of Lemma~\ref{har:lem:null-divisors} for
every \(N_i\), and hence \(N_0\), forces the intersection of \(f^*A\) with
\((\pi^*\theta)^r(\pi^*\omega)^{n-r-1}\) to vanish. Where the map to
the image has maximal rank this proves the kernel condition, by the
same positive-semidefinite linear algebra used above.

For every semigroup element \(N\) and integer \(e\ge0\),
\[
 \dim_{\mathbb C(Y)}H^0(V_\eta,\pi^*N+eE)^T\le1.
\]
If two generic sections were independent, twist their invariant direct
image by a sufficiently ample base line, obtain global invariant
sections, and absorb that twist with \eqref{har:eq:absorb}. Since
\(\pi_*\mathcal O_V(eE)=\mathcal O_X\), this gives invariant sections
of a semigroup element whose ratio is not in \(\mathbb C(Y)\),
contradicting its construction.

Let \(F_i\) be the horizontal divisor of \(\pi^*s_i\). For
\(a_i<a_j<a_h\), positive integers \(u,v\) with
\(ua_i+va_h=(u+v)a_j\) yield two nonzero invariant generic sections
of the same line. Their ratio is constant on the generic fiber, so
\(uF_i+vF_h=(u+v)F_j\). Consequently \(F_i\) is affine in \(a_i\).
Its slope is effective because \(F_i\ge0\) and \(a_i\to\infty\).
The ratio \(s=\pi^*s_2/\pi^*s_1\), of degree \(b=a_2-a_1>0\),
therefore has no horizontal poles. In fact
\[
 \dim_{\mathbb C(Y)}H^0(V_\eta,kbD_V+eE)^T=1
 \quad(k,e\ge0).
\]
The displayed section \(s^ks_E^e\) gives the lower bound. For the
upper bound choose \(i\) large enough that
\(F_i-k(F_2-F_1)\ge0\). Multiplication by \(\pi^*s_i/s^k\) embeds
the space into the preceding rank-at-most-one space for \(N_i+eE\).
If the base is a point, \(s\) is already globally regular. Otherwise
Propositions~\ref{har:prop:maximum} and \ref{har:prop:euler} finish this
route. The stronger all-\(e\) generic-rank statement is retained here;
only \(e=1\) is needed for the metric and transfer.

\section{Extremal weights and a dominated base}\label{har:sec:extremal}

This route uses a source component for a one-parameter action rather
than the invariant index. It applies to a smooth projective rationally
connected \(S\) with smoothly semipositive \(L=-K_S\). The final
product argument will give invariance under an entire algebraic torus.
The source-character theorem of Section~\ref{inj:extremal-weights}
already gives this conclusion under weaker hypotheses. Here the
purpose is to obtain the character by harmonic vanishing: we construct
a base dominated by \(L\), then retract that base to its source so
that Euler interpolation takes place in one specified weight.

\begin{lemma}[Source weights and retraction]\label{har:lem:source}
Let \(S\) be smooth connected projective with a \(\mathbb C^*\)-action,
allowing the trivial action. Let \(S_0\) be its attracting component
for \(t\to0\), and write \(w_C(U)\) for the fiber weight of a
linearized line \(U\) on a fixed component \(C\).
A nonzero eigensection of \(U\) has weight at most \(w_{S_0}(U)\).
If \(U\) is nef, then \(w_C(U)\le w_{S_0}(U)\) for every fixed
component, and the multiplicity of weight \(mw_{S_0}(U)\) in the
virtual character \(\chi(S,mU)\) is
\(\chi(S_0,mU|_{S_0})\), for \(m>0\).
If \(U\) is nef, on a smooth equivariant resolution \(d:\widehat S\to S\) of the
source retraction \(p:\widehat S\to S_0\), there is an equivariant
nonzero bundle map \(p^*(U|_{S_0})\to d^*U\).
\end{lemma}

\begin{proof}
An ample line has a linearized positive power for a projective $\C^*$-action. Indeed the action's map to the Picard variety is constant, since a torus has no nonconstant homomorphism to an abelian variety. The pairs consisting of an element and an isomorphism from the pulled-back line to the line form an algebraic extension of $\C^*$ by the scalar $\C^*$; this extension is a torus and splits. This supplies the needed linearization. The section-weight bound, the comparison of fixed weights, and the index formula are Lemma~\ref{lem:source-character} for $T=\C^*$ and $U$ as given. The source retraction itself is the attracting-cell map of \cite{BB1973}. The additional assertion needed by the harmonic construction is its bundle map, which we now prove.

For the bundle map, identify the generic fibers of \(U\) and its
restriction at the source by transporting a vector along an orbit,
dividing by \(t^{w_{S_0}(U)}\), and taking the limit. The orbit
extension over \(\mathbb A^1\), over the relevant function field,
pulls the line back to an equivariantly trivial line of that weight.
This constructs a multiplicative rational identification, hence the
claimed rational map. If \(U\) is ample, a generated power has enough
source-weight sections to generate its restriction to \(S_0\). Their
pulled-back restrictions map to the original sections by the limit
construction, proving absence of poles. For nef \(U\), a negative
order in the rational map would give a negative order for
\(kU+H\) for large \(k\), contrary to ampleness. This proves regularity
along every prime divisor and hence everywhere by normality. The
trivial action has \(S_0=S\) and the identity map.
\end{proof}

\begin{proposition}[Attainment of the source weight]\label{har:prop:extremal}
Let \(S\) be smooth connected projective rationally connected, with
smoothly semipositive \(L=-K_S\), and fix a \(\mathbb C^*\)-action with
the natural anticanonical linearization. If \(w=w_{S_0}(L)\), then
some \(\ell>0\) has a nonzero section of \(\ell L\) of weight \(\ell w\).
\end{proposition}

\begin{proof}
\smallskip\noindent\textit{Determinant sections with a prescribed weight.}
A point is immediate. Otherwise \(S_0\) is rationally connected,
being a smooth rational image of \(S\). Thus
\(\chi(S_0,\mathcal O_{S_0})=1\). The highest-weight index in
Lemma~\ref{har:lem:source} is a nonzero polynomial in \(m\), so for
unbounded \(m\) a fixed group \(H^q(S,mL)\) contains weight \(mw\).
The smooth coefficient Lefschetz map is equivariant for a
circle-invariant Kähler form. It gives twisted sections
\[
 v_m\in H^0(S,\Omega_S^{\dim S-q}\otimes(m+1)L)
      =H^0(S,\bigwedge^q T_S\otimes mL)
\]
of weight \(mw\). For \(q=0\) they are the desired sections. Otherwise
saturate their generic span in \(\bigwedge^qT_S\), of rank \(e>0\),
and call it \(\mathcal U\). The determinant extraction used earlier
gives nonzero eigensections
\begin{equation}\label{har:eq:detweights}
 d_i\in H^0(S,\det\mathcal U+k_iL),\qquad
 \operatorname{wt}(d_i)=k_iw,\qquad k_i\to\infty.
\end{equation}
If \(D_i\) is the zero divisor, we claim
\[
 eL-c_1(\mathcal U)\text{ pseudoeffective},\qquad
 (e+k_i)L-D_i\text{ pseudoeffective}.
\]
For $q=\dim S$, saturation gives $\mathcal U=\bigwedge^{\dim S}T_S=L$ and $e=1$, so $eL-c_1(\mathcal U)=0$ directly. For $0<q<\dim S$, the cotangent-subsheaf theorem \cite[Theorem~4.1]{LMPTX2023} applies: \(\mathcal U\otimes(-L)\subset\Omega_S^{\dim S-q}\), and
exterior powers are direct summands of tensor powers in characteristic
zero; saturating can only add an effective determinant correction
in the required direction.

\smallskip\noindent\textit{A base on which dominated systems become rigid.}
Call a divisor class dominated if a positive multiple of \(L\) minus
it is pseudoeffective. Maximize the image dimension of linear systems
of dominated effective integral divisors, and choose a full system
\(|P_0|\) realizing this maximum. Rational connectedness gives
\(\operatorname{Pic}^0(S)=0\), so the system may be linearized and
taken equivariantly. Resolve its graph, take Stein factorization,
flatten by the Hilbert-family construction, resolve the base and
normalize the main component. We obtain
\[
 V\longrightarrow W\longrightarrow S,\qquad g:W\to Y,
\]
with smooth projective \(V,Y\), normal equidimensional \(W\), and
connected fibers. Write \(f:V\to Y\), \(\mu:V\to S\),
\(A=\mu^*L\), \(E=K_{V/S}\), and let \(L_W\) be the pullback of \(L\) to \(W\). An ample pullback from \(Y\) is
dominated: the original moving polarization is dominated, and its
pullback to a birational base model is big, hence dominates a small
ample class. Also \(E\) is dominated. Horizontal exceptional primes
have centers in the original base locus and are bounded by pullbacks
of divisors in \(|P_0|\); vertical primes are bounded by sufficiently
ample base pullbacks. These statements prove domination of the whole
finite effective canonical discrepancy.

A line \(B\) on \(V\) is dominated when \(cA-B\) is pseudoeffective
for some \(c>0\). No such line can have two independent sections on
the generic fiber. If it did, tensor by a sufficiently ample base
pullback to globalize them. This new line remains dominated because
the base polarization is dominated. Combine their ratio with the
original base coordinates, using products to put all coordinates
in a single linear system. Its line is still dominated. Push the
effective system and its pseudoeffective domination relation down
to \(S\). The new ratio is nonconstant along the generic fiber, so
the image dimension exceeds the chosen maximum, a contradiction.
In particular generic-fiber sections of both \(jaA\) and
\(E+jaA\) have dimension at most one for every \(a,j>0\), and the
same adjoint assertion for \(j=0\) follows from domination of \(E\).

There are only finitely many dominated horizontal primes on \(V\).
Here \(\operatorname{Pic}^0(V)=0\), since \(V\) is birational to the
rationally connected \(S\). Thus its Picard group has finite rank.
Infinitely many dominated horizontal primes would give, after clearing
torsion, a nontrivial integral linear equivalence among them. Splitting
positive and negative coefficients gives two effective divisors in
the same class and a pencil varying along the generic fiber, contrary
to the preceding paragraph. The pullbacks of \(D_i\) therefore have
horizontal supports in a fixed finite set. Order their multiplicities
along a subsequence. A ratio from \eqref{har:eq:detweights} gives a
rational eigensection \(s\) of \(aL_W\), weight \(aw\), with no
horizontal poles.

\smallskip\noindent\textit{The metrized base line and its source weight.}
Divisorial normalization, as in
Proposition~\ref{har:prop:maximum}, gives
\begin{equation}\label{har:eq:weightbase}
 g^*Q\longrightarrow aL_W,
\end{equation}
with effective zero divisor omitting one lift over every base prime.
Shift the linearization of \(Q\) to make this map equivariant. Since
the canonical exceptional section times \(s^j\) is nonzero generically,
all required adjoint ranks are exactly one. The maximum construction
gives a continuous semipositive metric on \(Q\). Here the base action
may be nontrivial, but full rank one is available and the ordinary
direct-image line theorem applies without taking an invariant summand.

We must still control the character produced by interpolation on the
base. Let \(Y_0\) be the source of its induced action. The original
unbounded eigensections imply
\begin{equation}\label{har:eq:weightinequality}
 w_{Y_0}(Q)\ge aw.
\end{equation}
After fixing \(a\), take an infinite subsequence of \(k_i\) in one
residue class modulo \(a\), still ordered horizontally. Divide by its
first section. The resulting eigensections of \(jaL_W\), for
unbounded \(j\), have weight \(jaw\), no horizontal poles and vertical
poles bounded by one fixed divisor. Clear these poles by an invariant
effective Cartier divisor \(C\) on \(Y\), with its canonical section
of weight zero. Generic rank one and the map
\eqref{har:eq:weightbase} show that each resulting section comes from
\(jQ+C\). It has no residual poles: at a base prime test on a lift
omitted by the zero divisor of \eqref{har:eq:weightbase}. The source
upper bound gives
\(jaw\le j w_{Y_0}(Q)+w_{Y_0}(C)\). Divide by \(j\) and pass to the
unbounded subsequence to obtain \eqref{har:eq:weightinequality}.

\smallskip\noindent\textit{Euler interpolation on the source of the base.}
The continuous semipositive metric makes \(Q\) nef by
Lemma~\ref{har:lem:smoothing}. Resolve the source retraction of \(Y\)
and dominate this resolution by a smooth model \(V'\) of \(V\).
Let \(p:V'\to Y_0\), \(\mu':V'\to S\), \(A'=\mu'^*L\),
\(E'=K_{V'/S}\), and \(Q_0=Q|_{Y_0}\). Lemma~\ref{har:lem:source}
and \eqref{har:eq:weightbase} give
\(p^*Q_0\to aA'\). An ample pullback by \(p\) is dominated by
\(A'\): the pullback of the original ample base polarization is big
on the resolution of \(Y\), and hence dominates an ample class there,
which in turn dominates an ample pullback from \(Y_0\).
Proposition~\ref{har:prop:dominationvanish} applies to every twist
\(jQ_0\), including zero. The action on \(Y_0\) is trivial, and
\(Q_0\) has the constant fiber weight \(w_{Y_0}(Q)\). Let
\(\mathcal B\) be the weight-zero summand of \(p_*\mathcal O(E')\).
Thus sections of \(\mathcal B\otimes Q_0^j\) have precisely weight
\(j w_{Y_0}(Q)\). Vanishing identifies their dimension with the
Euler polynomial for every \(j\ge0\). At zero this polynomial is
positive, because the canonical exceptional section is invariant.
For some \(j>0\) we therefore obtain a nonzero section of
\(E'+jp^*Q_0\) of that weight. Map it into \(E'+jaA'\) and push
down to \(S\). Its weight is at least \(jaw\) by
\eqref{har:eq:weightinequality}, and at most \(jaw\) by the source
bound on \(S\). Equality gives the required extremal section.
\end{proof}

\begin{corollary}[From extremal characters to invariants]\label{har:cor:torus}
Every algebraic torus acting on the rationally connected \(S\) of
Proposition~\ref{har:prop:extremal} has a nonzero invariant section of
some positive power of \(-K_S\).
\end{corollary}

\begin{proof}
For a generic integral one-parameter subgroup \(\lambda\), its source
\(C\) is a torus-fixed component with all normal weights positive
on \(\lambda\). Denote by \(\gamma_C\) the fiber character of \(L\).
Choose a full torus eigenvector in the extremal section space of
Proposition~\ref{har:prop:extremal}. It cannot vanish on \(C\), since
a nonzero leading conormal jet would strictly lower its
\(\lambda\)-weight. Its character is therefore \(\ell\gamma_C\).
The finitely many characters of such sources have convex hull
containing zero. Otherwise a separating functional, approximated by
a generic rational one, would be negative on all of them, whereas
the source for that one-parameter subgroup has
\(\gamma_C(\lambda)>0\), the sum of its positive normal weights.
Use the effective torus quotient if the action has a kernel; a trivial
action is already covered. A rational convex combination of the source
characters is zero. Clear its denominators and the finitely many
section exponents, then multiply the corresponding powers of their
nonzero eigensections. On the integral variety this product is nonzero
and invariant.
\end{proof}

\section{The curvature-kernel field and a reduced flat alteration}\label{har:sec:alteration}

This construction converts ordinary cohomology on $X$ itself; in the descent application, $X$ will be a finite \'etale cover. It does not require \(\operatorname{Pic}^0=0\).
Its alteration is genuinely generically finite. We will explicitly
descend the resulting metrized base line to birational models of the
original function field.

\begin{theorem}[Ordinary conversion through a curvature-null field]\label{thm:alteration-conversion}
Let $X$ be smooth connected projective, with smoothly semipositive $L=-K_X$. Suppose that for one fixed $q\ge0$, $H^q(X,mL)\ne0$ for unbounded positive integers $m$. Then $H^0(X,rL)\ne0$ for some $r>0$.
\end{theorem}

Smooth Lefschetz and Corollary~\ref{inj:forms-input} already imply this
conclusion. We give a different geometric proof through three
constructions: the field constant in curvature-null directions, a line
of multiplication maps on a reduced flat alteration, and descent of
that metrized line and its injection to the original function field.
The alteration makes the fiberwise maximum positive and continuous by
using reduced fibers. Returning to the original function field then
allows the birational Euler argument to apply.

\subsection{The field constant in curvature-null directions}\label{har:subsec:kernelfield}

Let \(X\) be smooth connected projective with smoothly semipositive
\(L=-K_X\), and suppose \(H^q(X,mL)\ne0\) for arbitrarily large
positive \(m\), with fixed \(q\). Put \(n=\dim X\), let \(\theta\)
be the smooth semipositive curvature of \(L\), and let \(r\) be its
maximal rank. If \(q=0\), the conclusion is already given. Otherwise
start from the assumed classes and apply the harmonic-star construction
in the proof of Lemma~\ref{har:lem:forms}; it gives nonzero sections
of \(\Omega_X^{n-q}\otimes L^{m+1}\). No invariant-index assumption
is involved. If \(q=n\), these are anticanonical sections; if
\(r=n\), the positive curvature energy on an open set excludes
\(q>0\). In the remaining cases the same determinant extraction gives
a line \(M\) mapping nontrivially to a positive exterior power of
\(\Omega_X^{n-q}\), and effective divisors
\[
 D_i\in|M+a_iL|,\qquad a_i\to\infty,
 \qquad D_i[\theta]^r[\omega]^{n-r-1}=0.
\]
The intersection equality is Lemma~\ref{har:lem:null-divisors}.
If \(r=0\), all \(D_i\) vanish, and two distinct \(a_i\) show that
a positive power of \(L\) is trivial. Assume \(1\le r<n\).

On the ordinary open \(U\) of maximal curvature rank, define
\[
 \mathcal K=\{g\in\mathbb C(X): dg(\ker\theta)=0
       \text{ on }U\text{ off a proper algebraic subset}\}.
\]
It is a field. It is relatively algebraically closed: differentiate a
minimal polynomial for an element algebraic over \(\mathcal K\) in a
null direction; its separable derivative is generically nonzero, so
the element has zero derivative there as well. The intermediate-field
argument from Section~\ref{har:subsec:small} proves finite generation.

If a rational function has null pole support, it belongs to
\(\mathcal K\). Its level divisors have the same zero intersection
as the poles, hence each of their components is null. At a regular
level point in \(U\), the tangent hyperplane contains \(\ker\theta\),
by the same rank-quotient argument as before. This proves the assertion
where the differential is nonzero and then everywhere required.

Choose a smooth projective model \(Y\) of \(\mathcal K\) and resolve
the map by \(\mu:X_1\to X\), \(f:X_1\to Y\). Put \(A=\mu^*L\).
The generic fiber is geometrically integral, because the extension
is regular in characteristic zero. On a nonempty ordinary open,
\(\ker\mu^*\theta\subset\ker df\), and \(\dim Y\le r\).
In fact for every smooth \((1,1)\)-form \(\eta\) on \(Y\),
\begin{equation}\label{har:eq:globalwedge}
 (\mu^*\theta)^r\wedge f^*\eta=0.
\end{equation}
On the rank-\(r\) open this follows because every horizontal
covector lies in the positive-curvature span; elsewhere the rank is
smaller, and continuity extends the identity across the algebraic
exceptional locus.

Any integral relation \(\sum b_i=\sum b_i a_i=0\) makes
\(\sum b_iD_i\) principal with null pole support. Its function is in
\(\mathcal K\), so the horizontal parts of \(\mu^*D_i\) satisfy the
same relation. Three-term interpolation proves that these horizontal
divisors are affine in \(a_i\). Their slope is effective by unbounded
nonnegative coefficients. Thus for two indices with \(b=a_2-a_1>0\)
the quotient \(s\) is a rational section of \(bA\) with effective
horizontal divisor, regular over a dense base open. If \(Y\) is a
point it already descends to a section on \(X\).

Let \(E=K_{X_1/X}=K_{X_1}+A\). For every \(t\ge0\),
\begin{equation}\label{har:eq:ordinaryrank}
 \operatorname{rank}f_*\mathcal O_{X_1}(E+tA)\le1,
 \quad\text{with equality for }t=kb,\ k\ge0.
\end{equation}
For the upper bound, twist two hypothetical independent generic
sections by a high ample base line to make them global. Their zero
divisors push to effective divisors on \(X\) with zero intersection
against \([\theta]^r[\omega]^{n-r-1}\), using
\eqref{har:eq:globalwedge}, exceptionality of \(E\), and
\(\theta^{r+1}=0\). Their quotient has null poles, hence lies in
\(\mathcal K\), contrary to independence. The lower bound is supplied
by \(s_Es^k\).

Over a smooth dense open with geometrically integral fibers, the
function \(-\log\max_{X_{1,y}}|s|^2\) is plurisubharmonic. Indeed the
rank-one adjoint section \((s_E/f^*\eta)s^k\) has smooth coefficient
\((kb+1)A\). The moment calculation
\eqref{har:eq:moments} applies, with full-support density and locally
uniform power-mean limit. The next construction extends this good-locus
metric using reduced flat fibers rather than divisorial minima.

\subsection{A Hom line on a reduced flat family}\label{har:subsec:homline}

By weak semistable reduction \cite[Theorem~0.3]{AK2000},
there are a smooth projective alteration \(Y_a\to Y\) and a
modification \(Z_a\) of the main component of \(X_1\times_Y Y_a\)
such that \(p_a:Z_a\to Y_a\) is equidimensional, toroidal, and has
reduced fibers. Toroidal varieties are Cohen--Macaulay, so the
equidimensional flatness criterion over the smooth base makes
\(p_a\) flat. Write \(q_a:Z_a\to X_1\). The source \(Z_a\) need not
be smooth.

Take a sufficiently high power \(\mathcal P\) of a relatively very
ample line. Relative Serre vanishing and base change give locally
free bundles
\[
 \mathcal H=p_{a*}\mathcal P,\qquad
 \mathcal H'=p_{a*}(\mathcal P\otimes q_a^*A^b),
\]
commuting with base change, with surjective evaluation for
\(\mathcal P\). On the generic fiber, multiplication by \(s\) gives
a nonzero map \(\mathcal H_\eta\to\mathcal H'_\eta\). Resolve its
rational point in the projective bundle of lines in
\(\operatorname{Hom}(\mathcal H,\mathcal H')\) on a smooth projective
birational base model \(Y_b\to Y_a\), and call the pulled-back
tautological subline \(B_b\). Put \(Z_b=Z_a\times_{Y_a}Y_b\), with
maps \(p_b\) and \(q_b\). Flatness and reduced geometric fibers
survive this base change. The total space is integral: flatness over
an integral base embeds its local rings in their generic localizations,
and its generic fiber is integral.

There is a regular line map
\begin{equation}\label{har:eq:homline}
 p_b^*B_b\longrightarrow q_b^*A^b
\end{equation}
nonzero on each fiber. To construct it, compose the tautological Hom
map with evaluation into \(\mathcal P\otimes q_b^*A^b\). It
annihilates the kernel of evaluation onto \(\mathcal P\), because
it is multiplication on the generic fiber and the total space is
integral. Factor through \(\mathcal P\) and cancel that invertible
line. If the result vanished as a map on a fiber, the induced map
of fiber section spaces would vanish, contrary to the everywhere
subbundle \(B_b\subset\operatorname{Hom}(\mathcal H,\mathcal H')\)
and base change. Reducedness now implies that it has a nonzero
value at some point of every fiber; without reducedness this last
implication would fail.

Give \(B_b\) the maximum metric through \eqref{har:eq:homline} and
the pulled-back smooth metric of \(A^b\). The maximum is positive
and continuous. Properness gives upper semicontinuity. Smooth points
are dense in each reduced fiber; a value arbitrarily close to the
maximum may be chosen at such a point. Flatness and fiber smoothness
make the morphism smooth there, giving liftings of nearby base points
and lower semicontinuity. On a dense open the map is the pullback of
\(s\) times a rational base factor, by its defining generic Hom line.
Its fibers map onto the original smooth integral fibers, so the
maximum weight equals the pullback of the good-locus weight just
proved, plus a pluriharmonic frame term. Continuity and removable
singularities extend plurisubharmonicity everywhere on \(Y_b\).

\subsection{Descending the alteration and its section}\label{har:subsec:alterationnorm}

Proposition~\ref{har:prop:euler} requires a source birational to \(X\),
whereas the construction above used an alteration. We first descend
a power of the metrized line by a finite group quotient. We then
descend its injection by a field norm. These two operations will
produce the required data on birational models of the original fields.

Choose a finite Galois extension of \(\mathbb C(Y)\) containing
\(\mathbb C(Y_b)\), of degree \(\ell\) and group \(G\). Starting
with the normalization of \(Y\) in this extension, normalize the simultaneous graph of
all conjugate rational maps to \(Y_b\). This gives a normal projective
\(W\), a regular \(G\)-action, and a map \(\rho:W\to Y_b\).
The quotient \(W_0=W/G\) is normal projective birational to \(Y\),
and \(W\to W_0\) is finite. Let \(B_W=\rho^*B_b\). The product
\[
 C_W=\bigotimes_{g\in G}g^*B_W
\]
has its permutation linearization and invariant product metric.
The power \(C_W^\ell\) has trivial stabilizer actions on fibers,
so descends to a line \(C_0\) on \(W_0\). One can check descent
locally on affine quotient neighborhoods: prescribe an equivariant
nonzero value on a finite orbit, lift it to a section, average, and
shrink the quotient neighborhood until it is a frame. Its invariant
transition functions descend. The invariant metric descends
continuously by the quotient topology.

Resolve \(\widehat Y\to W_0\), and let \(B_0\) be the pullback of
\(C_0\). Its metric is plurisubharmonic on the dense finite-étale
locus, and by continuity everywhere. Choose a smooth projective
birational model \(\widehat X\to X_1\) mapping to \(\widehat Y\),
with composite maps $\widehat\mu:\widehat X\to X$ and $\widehat f:\widehat X\to\widehat Y$, and put \(\widehat A=\widehat\mu^*L\). Normalize the main component of
\(\widehat X\times_{W_0}W\), calling it \(R\). It is finite over
\(\widehat X\) with the same Galois group: geometric integrality of
the generic fiber, proved above, makes the function field extension
linearly disjoint from the finite Galois base extension.
The finite maps fit into the commutative square
\[
\begin{array}{ccc}
 R & \longrightarrow & W\\
 \mathllap{\scriptstyle\mathrm{finite}\,}\big\downarrow &&
       \big\downarrow\mathrlap{\,\scriptstyle\mathrm{finite}}\\
 \widehat X & \longrightarrow & W_0 .
\end{array}
\]
Here the lower map factors through \(\widehat f:\widehat X\to
\widehat Y\), and the upper-right space maps to \(Y_b\) by \(\rho\).

On \(R\), the difference of the pullbacks \(b\widehat A-B_W\) has
a nonzero section. Indeed resolve the rational map from \(R\) to
\(Z_b\) over \(X_1\) and \(Y_b\). The resulting common graph is a
modification of \(R\), and \eqref{har:eq:homline} gives a regular
section there. Both line bundles already come from \(R\); normality
therefore descends this section to \(R\). Multiply its \(G\)-conjugates and
raise to the \(\ell\)-th power. This yields a section of the
pullback of \(b\ell^2\widehat A-\widehat f^*B_0\). A further field
norm for the finite surjection onto the normal target
\cite[Lemma~31.18.7]{StacksNormNormal}, of degree \(\ell\), relative to a rational line frame, gives
a regular nonzero section downstairs of
\begin{equation}\label{har:eq:descendedalteration}
 b\ell^3\widehat A-\widehat f^*(\ell B_0).
\end{equation}
Regularity can be checked at every prime: all orders upstairs are
nonnegative, hence so are the norm orders. The powers are not meant
to be minimal. Put \(\widehat b=b\ell^3\) and
\(\widehat Q=\ell B_0\). We have returned to birational models of
the original source and base fields, with a continuously semipositive
\(\widehat Q\) and an injection
\(\widehat f^*\widehat Q\to\widehat b\widehat A\).
The original kernel condition persists on a nonempty open of these
birational models. Theorem~\ref{har:thm:kernel} and
Proposition~\ref{har:prop:euler} now give a section of a positive
power of \(L\) on \(X\).

\section{Local strictness and an exact multiplier ideal}
\label{inj:localized}

The rank-one direct image of an adjoint line need not be locally free.
We identify its defect from its reflexive hull with the multiplier ideal
of its fiber-integral metric. This exact identity permits Euler
interpolation without discarding conditions in codimension two.

The two analytic lemmas below turn positivity on one horizontal patch
into vanishing for that ideal. Proposition~\ref{inj:exact-adjoint-ideal}
then identifies the ideal, and Proposition~\ref{inj:degree-zero-route}
applies the construction to the curvature-null-pole field already used
in Section~\ref{har:subsec:small}. The base field is the same; the
vanishing argument now uses the actual degree-zero direct image.

\begin{lemma}[Localized multiplier-ideal vanishing]\label{inj:local-nadel}
Let $Y$ be a smooth connected projective complex variety of dimension
$d>0$.  Suppose a line bundle $R$ has a singular Hermitian metric with
plurisubharmonic weights $\psi$.  Assume that on some coordinate ball
these weights are continuous and their curvature dominates a positive
smooth $(1,1)$-form.  Then
\[
 H^q\bigl(Y,\mathcal O_Y(K_Y+R)\otimes\mathcal I(\psi)\bigr)=0
 \qquad(q>0).
\]
\end{lemma}

\begin{proof}
We construct a globally strictly positive metric with exactly the same
multiplier ideal.

\smallskip\noindent\textit{Producing an integrable divisor metric.}
Choose a very ample $A$ with curvature $\omega_A$,
and a smooth metric $\psi_0$ on $R$.  For sufficiently small fixed
$\delta>0$ and arbitrarily large integers $l$, we first claim that
\begin{equation}\label{inj:jet-section}
 H^0\bigl(Y,\mathcal O_Y(lR-A)\otimes
       \mathcal I(l(1+\delta)\psi)\bigr)\ne0.
\end{equation}
On $lR$ use weights $l(1+\delta)\psi-l\delta\psi_0$, and add the
smoothly metrized line $j_lA-K_Y$, with
$j_l=\lceil C\delta l+C'\rceil$.  Fixed sufficiently large constants
$C,C'$ absorb the negative smooth curvatures and leave a positive
lower bound.  The multiplier ideal is
$\mathcal I_l=\mathcal I(l(1+\delta)\psi)$.
Within a smaller ball of strict positivity choose a point and a cutoff
logarithmic pole equal to $\log|z|^2$ near that point.  Add $\eta l$
times this pole, with $\eta>0$ fixed sufficiently small independently
of $\delta,l$.  Its negative curvature costs at most $O(\eta l)$
on the ball, where the original positive lower bound has order $l$.
Nadel vanishing \cite{Nadel1990}, in the strict-current form of
\cite[Theorem~5.11]{Demailly2011}, therefore gives surjectivity to jets
of order $\lfloor\eta l\rfloor-d$ at the point.  Continuity of
$\psi$ there makes $\mathcal I_l$ trivial near that point, and
polar-coordinate integration gives the ideal
$\mathfrak m^{\lfloor\eta l\rfloor-d+1}$ for the added pole.
Consequently
\[
 h^0\bigl(Y,(lR+j_lA)\otimes\mathcal I_l\bigr)\ge c l^d
\]
for some $c>0$ independent of sufficiently small $\delta$.

Remove $j_l+1$ copies of $A$.  At each step choose a smooth member of
$|A|$ containing no associated center of $\mathcal O_Y/\mathcal I_l$.
If a section vanishes on that member, division by its equation
preserves membership in $\mathcal I_l$, because that equation is a
nonzerodivisor on the quotient.  The number of conditions lost at a
step is at most $C''l^{d-1}$ uniformly for $0<\delta\le1$.
Indeed choose $p$ with $pA-R$ globally generated, restrict a general
nonzero section to the hyperplane, and bound the restriction space by
that of $(pl+h)A$ there, where $0\le h\le j_l$.  The Hilbert
polynomial and the hyperplane exact sequence give this bound.
The total loss is $O(\delta l^d)+O(l^{d-1})$.  Choosing $\delta$
small and then $l$ large proves \eqref{inj:jet-section}.

Let $\sigma$ be a nonzero section in \eqref{inj:jet-section}. In compatible
local frames, the weight $\phi_A$ of the positive metric on $A$ defines a divisor
metric on $R$ by
\[
 \psi'=\frac{1}{l}\bigl(\log|\sigma|^2+\phi_A\bigr).
\]
Its curvature dominates $\omega_A/l$. If
$V=\psi'-(1+\delta)\psi$, membership of $\sigma$ in the multiplier ideal
in \eqref{inj:jet-section} says exactly that $\exp(lV)$ is locally
integrable.

\smallskip\noindent\textit{Preserving the ideal.} Take
$l\ge1/\delta$ and put
$\rho=(1-\varepsilon)\psi+\varepsilon\psi'$.  The identity
\[
 -\psi=-\frac{\rho}{1+\varepsilon\delta}
             +\frac{\varepsilon V}{1+\varepsilon\delta}
\]
and H\"older's inequality, with the integrability of
$\exp(V/\delta)$, prove
$\mathcal I(\rho)\subseteq\mathcal I(\psi)$.
For the converse, strong openness
\cite[Theorem 1.1]{GuanZhou2015}, coherence, and compactness give
$\gamma>0$ with
$\mathcal I((1+\gamma)\psi)=\mathcal I(\psi)$ on a finite cover.
Small-exponent integrability of the divisor metric $\psi'$ gives
$\gamma'>0$ with $\exp(-\gamma'\psi')$ locally integrable.
Apply H\"older with
$s=(1-\varepsilon)/(1+\gamma)$ and $1-s$.  For small
$\varepsilon$ one has $\varepsilon/(1-s)<\gamma'$, proving
$\mathcal I(\psi)\subseteq\mathcal I(\rho)$.
The metric $\rho$ has a global positive curvature lower bound and
the same ideal.  Nadel vanishing proves the assertion.
\end{proof}

\begin{lemma}[Strictness from horizontal curvature]
\label{inj:strict-direct-image}
Let $p:V\to Y$ be a proper holomorphic submersion with compact fibers
and K\"ahler total space.  Let $D$ have a smooth semipositive
Hermitian metric.  Work on an open set where
$p_*(K_{V/Y}+D)$ is locally free and its fibers are the corresponding
section spaces.  Suppose the coefficient curvature dominates the
pullback of a positive base form on an open patch meeting $V_{y_0}$,
and let $u_0\in H^0(V_{y_0},K_{V_{y_0}}+D|_{V_{y_0}})$ be a
section that is nonzero on that patch. The curvature of the
direct-image metric paired with $u_0$ at $y_0$ is strictly positive
in every nonzero tangent direction of $Y$.
If a compact torus acts over the trivial base action and preserves the
metrics, let the entire invariant summand, or an entire character
summand, have rank one and suppose $u_0$ belongs to that summand.
Its quotient metric under the holomorphic orthogonal projection has
the same strictness at $y_0$.
\end{lemma}

\begin{proof}
We spell out the consequence of Berndtsson's curvature formula
\cite[Section 4, Proposition 4.2, Lemmas 4.3--4.4, and equation (4.8)]{Berndtsson2009},
since strict vertical positivity is not required.  Restrict to a disk
with coordinate $z$ in a chosen base direction, and let the fiber
dimension be $a$. If $a=0$, the direct image is locally the orthogonal
sum of the coefficient lines on the finitely many sheets. Its
curvature pairing is a sum of nonnegative terms, with a positive
term on the patch where $u_0$ is nonzero. Assume now $a>0$.
A holomorphic direct-image section extending $u_0$
with vanishing covariant derivative at zero has a smooth absolute
$(a,0)$ representative $u$ satisfying
\[
 D'u=dz\wedge\lambda,\qquad
 \bar\partial u=dz\wedge\eta,
 \qquad D'=\partial-\partial\phi\wedge.
\]
On the central fiber, $\lambda$ is orthogonal to holomorphic adjoint
forms and $\eta\wedge\omega$ is $\bar\partial$-exact.
Replace $u$ by $u-dz\wedge\xi$ so that $\lambda=0$ and
$\eta$ is primitive there.  These two requirements prescribe,
respectively, $\bar\partial^*(\omega\wedge\xi)$ and
$\bar\partial(\omega\wedge\xi)$, up to the K\"ahler-identity
constants.  Orthogonality and exactness place the two prescriptions
in their required images; Hodge decomposition solves them
simultaneously.  Lefschetz linear algebra identifies the relevant
$(a,1)$-forms with $\omega\wedge\xi$, for $\xi$ of type
$(a-1,0)$.

With $c_a=i^{a^2}$, the resulting curvature pairing is
\[
 c_a p_*\bigl(u\wedge\bar u\,e^{-\phi}
                   \wedge i\partial\bar\partial\phi\bigr)
 -i\,dz\wedge d\bar z\,c_a
       \int_{V_0}\eta\wedge\bar\eta\,e^{-\phi}.
\]
The second term is nonnegative by the primitive middle-degree
Hodge--Riemann sign. The first integrand is nonnegative everywhere.
On the prescribed patch the pullback lower bound contributes the
positive base form times the squared vertical norm of $u_0$.
This is positive on a nonempty open subset, so its integral is
strictly positive. For a torus action, averaging
is a holomorphic orthogonal projection.  The Chern connection
preserves the orthogonal character splitting, so the same calculation
applies on the invariant summand.
\end{proof}

\begin{proposition}[The exact ideal of an invariant adjoint line]
\label{inj:exact-adjoint-ideal}
Let $X,V,Y$ be smooth connected projective varieties with $\dim Y>0$,
let $\pi:V\to X$ be birational, and let $f:V\to Y$ be surjective.
Suppose an algebraic torus $T$ acts equivariantly on these maps and
trivially on $Y$.
Give $L=-K_X$ its natural linearization and a smooth semipositive
metric invariant under the maximal compact torus. Put $L'=\pi^*L$,
$J=K_{V/X}$, and let $t$ be its natural invariant Jacobian section.
Assume that
\[
 \mathcal B=(f_*\mathcal O_V(J))^T
\]
has rank one. Via division by $t$, view $\mathcal B$ as a sheaf of
rational functions on $Y$, and write
$\mathcal B^{\vee\vee}=\mathcal O_Y(B_0)$ and $R=B_0-K_Y$.
Then the invariant adjoint $L^2$ metric extends to a singular
semipositive metric $\psi_R$ on $R$, and
\begin{equation}\label{inj:exact-ideal}
 \mathcal O_Y(B_0)\otimes\mathcal I(\psi_R)=\mathcal B.
\end{equation}
If, on a nonempty patch in the smooth base-change locus where $t$ is
nonzero, the curvature of $L'$ dominates the pullback of a positive
base form, then $\psi_R$ is smooth and strictly positive on a base
coordinate ball.
\end{proposition}

\begin{proof}
The sheaf $\mathcal B$ is torsion-free. Its hull is determined by
regularity over the generic points of base primes. More precisely,
for a prime $P\subset Y$ its coefficient in $B_0$ is
\[
 b_P=\min_{Q\mapsto P}
          \left\lfloor\frac{\operatorname{ord}_Q(t)}{
          \operatorname{ord}_Q(f^*P)}\right\rfloor.
\]
These are exactly the regularity tests for the rational section
$(f^*a)t$. On the smooth base-change open, the entire invariant
adjoint summand is a line. Its fiber-integral metric is semipositive
by the compact-torus orthogonal quotient argument of
Proposition~\ref{har:prop:maximum}.

Put $k=\dim Y$. This metric extends plurisubharmonically across every base prime.
At a general point of a component attaining $b_P$, write the map
as $(z_1^e,z_2,\ldots,z_k)$, where $e>0$ is the pullback
multiplicity, and let $a$ be the order of $t$.
The relative canonical coefficient for the local frame of $R$ is
a nonvanishing smooth factor times
\[
 z_1^{a-b_Pe-(e-1)}.
\]
Since $0\le a-b_Pe<e$, this exponent is nonpositive.
Integration on one fixed small vertical polydisk bounds the squared
norm below by a positive constant near the base prime.  The weight
is therefore locally bounded above there.  Extend first along these
divisorial neighborhoods and then across codimension at least two;
write the extended weight as $\psi_R$.

Choose local frames $g$ of $\mathcal O_Y(B_0)$ and $\eta$ of $K_Y$.
The frame $g/\eta$ of $R$ corresponds on the smooth locus to the
relative adjoint section $(f^*g)t/f^*\eta$. For a holomorphic
function $a$, Fubini identifies integrability of $a$ with this frame's
squared norm, using the base volume from $\eta$, with square
integrability of the absolute top form $(f^*(ag))t$ with the smooth
$L'$ metric upstairs. The excluded analytic sets have measure
zero.  Such an integrable meromorphic top form has no divisorial
pole; smoothness upstairs then extends it across codimension two.
Conversely a holomorphic lift is integrable over inverse images of
relatively compact base charts by properness.  The lift is invariant
and the identification agrees on the generic fiber.  This proves
both inclusions in \eqref{inj:exact-ideal}, so it retains any
codimension-two difference between $\mathcal B$ and its hull.

For the final assertion apply Lemma~\ref{inj:strict-direct-image}
to the invariant section determined by $t$ on the stated patch.
The fiber-integral metric is smooth on the base-change open, so
its strictness at one point gives a positive lower curvature bound
on a smaller coordinate ball.
\end{proof}

\begin{proposition}[Invariant sections by a degree-zero field]
\label{inj:degree-zero-route}
Let $X$ be smooth connected projective with
$H^q(X,\mathcal O_X)=0$ for $q>0$.  Suppose $L=-K_X$ has a smooth
semipositive metric and an algebraic torus $T$ acts on $X$.
Then $H^0(X,mL)^T\ne0$ for some $m>0$, with the natural
anticanonical linearization.
\end{proposition}

\begin{proof}
\smallskip\noindent\textit{The null-pole field and the relative section.}
Average a metric on $L$ and a K\"ahler form $\omega$ over the maximal
compact torus, and write $\alpha\ge0$ for the resulting curvature
form, in first-Chern normalization.  Proposition~\ref{har:prop:index}, with its actual value at zero,
followed by the equivariant hard Lefschetz construction in
Lemma~\ref{har:lem:forms},
has the following consequence: if the asserted invariant sections do
not exist, there is a fixed $q>0$ and unbounded positive $j$ with
\[
 0\ne u_j\in H^0(X,\Omega_X^{n-q}\otimes jL)^T.
\]
The coefficient in hard Lefschetz is $j=m+1$, since
$K_X+jL=mL$ equivariantly.  In the Bochner identity for the associated
harmonic $(n,q)$-forms the curvature term is a sum of $q$ eigenvalues
of $j\alpha$.  Full rank on an open set would force vanishing there,
hence vanishing of the holomorphic Lefschetz preimage everywhere.
Thus $r=\max\operatorname{rank}\alpha<n$.

Define the degree of a divisor by
\[
 \ell(D)=[D]\,[\alpha]^r[\omega]^{n-r-1}.
\]
It is nonnegative on effective divisors and vanishes on $L$.
Lemma~\ref{har:lem:forms} gives a linearized determinant line $S$
and, for unbounded positive integers $p$, nonzero invariant sections
\[
 s_p\in H^0(X,S+pL)^T.
\]
Put $N_p=\operatorname{div}(s_p)$. Lemma~\ref{har:lem:null-divisors}
gives $\ell(N_p)=0$: prescribed Ricci curvature first gives
$\ell(S)\le0$, and effectivity supplies the reverse inequality.
The degree-zero divisors are exactly the curvature-null divisors
of Section~\ref{har:subsec:small}. We use their invariant pole field
again, now to construct the exact direct-image ideal.

Let $K$ be the field of invariant rational functions whose polar
divisors have degree zero.  Sums and products preserve this property;
inverses do also, because polar and zero divisors are linearly
equivalent and effective degrees are nonnegative.  The field is
relatively algebraically closed in $\mathbb C(X)$.  Indeed a root
of a monic equation over $K$ can have poles only on the polar
divisors of its coefficients, by the discrete-valuation inequality.
The connected torus fixes each member of the finite root set, so the
root is invariant as well.  The field is finitely generated: choose
a transcendence basis over $\mathbb C$; algebraic subextensions
over its rational field inside the finitely generated ambient field
have bounded degrees, after adjoining an ambient separating
transcendence basis.  Their compositum is therefore finite.

Take a smooth projective model $Y$ of $K$.  Flatten the graph of the
induced rational map, and normalize its main component, to obtain a
normal equidimensional $G\to Y$ birational to $X$; then take an
equivariant resolution $V\to G$.  The action on $Y$ is trivial. Write $f_G:G\to Y$ for the normal-model morphism and $L_G$ for the pullback of $L$ to $G$.
Write $\pi:V\to X$, $f:V\to Y$, $L'=\pi^*L$, and let
$t$ be the invariant Jacobian section of
$K_V+L'=\mathcal O_V(J)$, with $J=K_{V/X}\ge0$.
The generic fiber is geometrically integral because $K$ is relatively
algebraically closed.  These constructions use only equivariant
resolution of the graph and the Hilbert-family flattening; they do
not assert equidimensionality of the smooth resolution.

Ratios of products of the $s_p$ having the same bidegree in $(S,L)$
belong to $K$.  Their horizontal divisor on $G$ is zero, so for
$p>p_1>p_0$,
\[
 (p-p_0)N_{p_1}^{\rm hor}
 =(p-p_1)N_{p_0}^{\rm hor}+(p_1-p_0)N_p^{\rm hor}.
\]
Equivalently,
\[
 N_p^{\rm hor}=N_{p_0}^{\rm hor}
 +\frac{p-p_0}{p_1-p_0}
       (N_{p_1}^{\rm hor}-N_{p_0}^{\rm hor}).
\]
Effectivity for unbounded $p$ forces
$N_{p_1}^{\rm hor}-N_{p_0}^{\rm hor}\ge0$.
Put $d_0=p_1-p_0>0$ and
$\sigma_0=s_{p_1}/s_{p_0}$, a rational invariant section of
$d_0L$ with effective horizontal divisor. We use the same symbol
for its pullback to $G$. If $Y$ is a point, it has no poles and
proves the assertion. Assume $\dim Y=k>0$.

For each prime $P\subset Y$, put
\[
 c_P=\min_{Q\mapsto P}
       \frac{\operatorname{ord}_Q(\sigma_0)}
            {\operatorname{ord}_Q(f_G^*P)}.
\]
Only finitely many $c_P$ are nonzero. Choose a positive integer $h$
clearing their denominators, and set
\[
 d=hd_0,\qquad \sigma=\sigma_0^h,\qquad
 B=\sum_P h c_P P.
\]
Thus $\sigma$ is a rational invariant section of $dL_G$, and $B$
is an integral divisor with trivial action. Divisorial normalization
gives
\begin{equation}\label{inj:invariant-base-decomposition}
 s_D=\sigma(f_G^*1_B)^{-1}
     \in H^0(G,dL_G-f_G^*B)^T,
 \qquad D=\operatorname{div}(s_D)\ge0.
\end{equation}
Here $1_B$ is the rational canonical section of $\mathcal O_Y(B)$.
At least one component above every base prime has coefficient zero
in $D$. We also write $\sigma$ for its pullback to $V$.

\smallskip\noindent\textit{Adjoint rigidity and a horizontal curvature bound.}
For every $m\ge0$,
\begin{equation}\label{inj:degree-rank}
 \operatorname{rank}
 \bigl(f_*\mathcal O_V(J+mdL')\bigr)^T=1.
\end{equation}
There is a generic section $t\sigma^m$.  If the invariant rank
were larger, a large very ample base twist would give two invariant
global sections independent over $K$.  Their pushed divisors on $X$
have degree zero.  For the base twist this follows by taking the
ratio of two general base hyperplanes: their pullbacks to the
equidimensional $G$ have no common prime component, so their pushed
divisors are the actual polar and zero divisors of an element of
$K$.  The remaining class is $mdL$, and the exceptional Jacobian
pushes to zero.  Thus the ratio of the two sections would again
belong to $K$, contradicting their independence.  This proves
\eqref{inj:degree-rank}, including $m=0$.

At a rank-$r$ point of $\alpha$, the differential of every rational
function in $K$ kills $\ker\alpha$ wherever it is regular.
For a regular value, the closure of its level hypersurface has
degree zero.  The nonnegative form
$\alpha^r\wedge\omega^{n-r-1}$ therefore vanishes on that
hypersurface.  A tangent hyperplane failing to contain $\ker\alpha$
would still carry rank $r$, a contradiction.  Apply this to $k$
independent base coordinates. It gives $k\le r$. Intersect the
maximal-rank open with the dense smooth base-change locus and the
complement of the exceptional locus. On a small patch there,
\begin{equation}\label{inj:horizontal-patch}
 \pi^*\alpha\ge c f^*\omega_Y\qquad(c>0).
\end{equation}
This is semipositive linear algebra on a complement of the kernel,
not an inference from numerical degree alone.  If $r=0$, a nonzero
effective divisor cannot have degree zero, so $k>0$ cannot occur.

Apply Proposition~\ref{har:prop:maximum} to the equivariant normal equidimensional model $G\to Y$, with $P=L$ naturally linearized and the rational section $\sigma$. Equation~\eqref{inj:degree-rank} verifies the invariant adjoint ranks, and \eqref{inj:invariant-base-decomposition} is its vertical normalization. The result is a positive continuous semipositive metric on $B$.

\smallskip\noindent\textit{The exact ideal and interpolation from zero.}
The varying twist is now metrized. The degree-zero metric must retain
the entire direct image, including its codimension-two defect.

Apply Proposition~\ref{inj:exact-adjoint-ideal} to
$\mathcal B=(f_*\mathcal O_V(J))^T$. Equation~\eqref{inj:degree-rank}
at $m=0$ gives its rank, and the patch in
\eqref{inj:horizontal-patch} is disjoint from the exceptional locus,
so $t$ is nonzero there. The proposition supplies the metric
$\psi_R$ on $R=B_0-K_Y$, strictly positive on a base ball, with the
exact identity \eqref{inj:exact-ideal}.

Adding $m$ times the continuous bounded metric on $B$ preserves the
ideal and the strictness on the ball.  Lemma~\ref{inj:local-nadel}
therefore gives
\[
 H^q(Y,\mathcal B\otimes\mathcal O_Y(mB))=0
 \qquad(q>0,\ m\ge0).
\]
At zero, canonical pushforward gives
$h^0(Y,\mathcal B)=h^0(V,\mathcal O_V(J))^T=1$.
Thus the Euler polynomial has value one at zero.  For some $m>0$ lift a nonzero section
and multiply by $s_D^m$.  The result is an invariant section of
$J+mdL'$.  Remove the exceptional Jacobian divisor by pushforward
to the smooth $X$.  It gives the required natural invariant section
of $mdL$.
\end{proof}

\section{Actual direct-image lines and stabilized orders}
\label{inj:stabilized}

This proof of natural invariant transfer retains the actual adjoint
direct-image sheaves throughout. A birational toroidal equidimensional
model makes their invariant rank-one summands line bundles. Local
frames of these lines lift to holomorphic adjoint sections, which will
give the integrability needed for their quotient metrics.

The boundary orders are minima of finitely many affine functions.
After the slopes are made integral, the lines vary affinely in all
large degrees. The intercept comparison also embeds the zero-degree
line, twisted by the eventual slope, into those large-degree lines.
Strong openness then carries strict positivity from degree zero along
this family, preserving the nonzero constant term in Euler interpolation.

\begin{proposition}[Invariant Iitaka conversion by stabilized orders]
\label{inj:stabilized-transfer}
Let $X$ be smooth connected projective and let an algebraic torus
$T$ act on $X$.  Give $L=-K_X$ its natural linearization and suppose
it admits a smooth semipositive metric.  Let $M$ be a linearized
line bundle with $-M$ pseudoeffective.  If there are nonzero
sections $\sigma_j\in H^0(X,M+jL)^T$ for an unbounded set
$J\subset\mathbb Z_{>0}$, then $H^0(X,mL)^T\ne0$ for some
integer $m>0$.
\end{proposition}

\begin{proof}
\smallskip\noindent\textit{Generic adjoint rigidity.}
Write $D(j)=M+jL$ and fix $j_0<j_1$ in $J$.
Choose $k>0$ so that the invariant system
$H^0(X,kD(j_1))^T$ has maximal image dimension among these systems
for all positive degrees.  No finite-generation assertion is needed
to attain this maximum, an integer between zero and $\dim X$.
An equivariant principalization of its base ideal gives a smooth
projective $\pi_0:W_0\to X$, isomorphic off the base locus, and
a connected-fiber Stein map $f_0:W_0\to Y_0$ such that
\[
 \pi_0^*kD(j_1)=f_0^*H+G,
\]
where $H$ is ample on the normal projective $Y_0$ and $G$ is the
effective fixed divisor.  The base action and the fiber action on
$H$ are trivial, while the defining section of $G$ is invariant.
These identifications follow by dividing the invariant generating
sections by the common fixed factor.  The action on the finite
Stein factor is trivial by connectedness of $T$.

Let $K=\mathbb C(Y_0)$ and let $F$ be the generic fiber.  It is
normal and integral, with $H^0(F,\mathcal O_F)=K$.
The field $K$ is relatively algebraically closed in $K(F)$:
an algebraic element has no poles on proper normal $F$, as follows
from a monic equation over $K$.  For every $h>0$,
\begin{equation}\label{inj:invariant-fiber-rigidity}
 \dim_K H^0(F,h\pi_0^*D(j_1)|_F)^T=1.
\end{equation}
There is the section $\pi_0^*\sigma_{j_1}^h$.  If there were
two independent invariants, their ratio would be transcendental
over $K$.  They extend as invariant rational sections on $W_0$
with only vertical poles.  A common pullback of a section of a
sufficiently large multiple $lH$ clears these poles: prescribe
vanishing on their finitely many proper base images.  Multiplying
also by $s_G^l$ produces invariant sections of a multiple of
$D(j_1)$ on $X$.  Products with the original system retain its
image field and add their new ratio, contradicting maximal image
dimension.  This proves \eqref{inj:invariant-fiber-rigidity}.

Let $D_i^F$ be the effective divisor of $\sigma_i|_F$.
The invariant products of the same linearized degree are proportional,
so for $c>j_1$ in $J$,
\[
 (c-j_0)D_{j_1}^F
 =(c-j_1)D_{j_0}^F+(j_1-j_0)D_c^F.
\]
Thus $D_{j_1}^F-D_{j_0}^F\ge0$.  Put $d=j_1-j_0>0$ and
$z=\sigma_{j_1}/\sigma_{j_0}$, a rational invariant section
of $dL$.  It is regular on $F$, with zero support in $D_{j_1}^F$.
If the base is a point we already have a regular invariant section
on $W_0$, hence on $X$.

Otherwise let $e_0$ be the Jacobian section of
$K_{W_0}+\pi_0^*L$.  Its zeros lie over the principalized base
locus and, on $F$, in the support of $D_{j_1}^F$.  For every fixed $s\ge0$, choose $h$ large enough that, on $F$,
\[
 hD_{j_1}^F\ge \operatorname{div}(e_0|_F)
                         +s\operatorname{div}(z|_F).
\]
Such an $h$ exists because both effective divisors on the right
have support in $D_{j_1}^F$.  The invariant quotient
$(\sigma_{j_1}|_F)^h/(e_0z^s|_F)$ is regular.  Multiplication
by this nonzero quotient injects the entire invariant adjoint
section space into the one-dimensional space in
\eqref{inj:invariant-fiber-rigidity}.  Consequently
\begin{equation}\label{inj:adjoint-generic-line}
 H^0\bigl(F,(K_{W_0}+(1+sd)\pi_0^*L)|_F\bigr)^T
       =K\cdot(e_0z^s)|_F\qquad(s\ge0).
\end{equation}

\smallskip\noindent\textit{Actual lines and affine boundary orders.}
The generic rank computation will now be upgraded to line bundles
whose local frames lift to adjoint sections. After further birational
modifications, the same adjoint ranks are preserved by canonical pushforward on a common smooth resolution.
No support claim about new exceptional primes is needed.
Now apply birational weak toroidalization
\cite[Theorem 1.1]{ADK2013}, followed by equidimensionalization
by toroidal subdivisions \cite[Proposition 4.4]{AK2000}.
This gives a birational modification $Y\to Y_0$ with $Y$ smooth,
a normal toroidal $Z$ birational over $W_0$, and an equidimensional
$g:Z\to Y$.  These are birational modifications; no alteration
or assertion of reduced fibers is being used.  Take a smooth
resolution $W\to Z$ and write $f:W\to Y$, $\pi:W\to X$,
and $\pi_Z:Z\to X$.

Toroidal singularities are rational.  Therefore, for
\[
 \mathscr E_s=f_*\mathcal O_W(K_{W/Y}+(1+sd)\pi^*L),
\]
the sheaf $\mathscr E_s\otimes\mathcal O_Y(K_Y)$ equals
\[
 g_*\bigl(\omega_Z\otimes
                  \pi_Z^*\mathcal O_X((1+sd)L)\bigr).
\]
It is reflexive: deleting a codimension-two subset of $Y$ deletes
only codimension at least two from $Z$, by equidimensionality,
and the reflexive canonical sheaf on $Z$ extends its sections.
The pushforward is torsion-free as well, since the source is
integral and the map is dominant.

The toroidal modifications may fail to be equivariant.
Nevertheless $\mathscr E_s$ has a $T$-linearization.  Compute the
same direct image on an equivariant smooth resolution $W'$ of
the main component of $W_0\times_{Y_0}Y$.  A common smooth
resolution of $W$ and $W'$ identifies their canonical direct
images, since all coefficient twists are pulled back from $X$
and a birational morphism of smooth varieties pushes its canonical
sheaf to the canonical sheaf.  The identifications are the
identifications of rational forms, so preserve the generic torus
action.  The invariant part
\[
 A_s=(\mathscr E_s)^T
\]
is now an actual reflexive direct summand of rank one by
\eqref{inj:adjoint-generic-line}; on smooth $Y$ it is a line
bundle.  Birational canonical pushforward gives
\begin{equation}\label{inj:global-invariant-identification}
 H^0(Y,K_Y+A_s)=H^0(X,sdL)^T,
 \qquad h^0(Y,K_Y+A_0)=1.
\end{equation}
The last equality uses the natural invariant trivialization of
$K_X+L$.

We compute these actual lines.  On the regular locus of $Z$ let
$e_Z$ be the Jacobian section with effective divisor $E_Z$.
Use the rational generator $e_Zz^s$ for $K_Y+A_s$.
For each prime $B\subset Y$, let $Q_i$ run over primes of $Z$
mapping onto $B$, and let
\[
 b_i=\operatorname{ord}_{Q_i}(g^*B)>0,
 \quad e_i=\operatorname{ord}_{Q_i}(E_Z),
 \quad c_i=\operatorname{ord}_{Q_i}(z).
\]
A rational base multiplier $a$ is regular above the generic point
of $B$ precisely when
$b_i\operatorname{ord}_B(a)+e_i+sc_i\ge0$ for all $i$.
There are no horizontal poles by
\eqref{inj:adjoint-generic-line}.  The exact allowed pole order is
\[
 u_s(B)=\left\lfloor\min_i\frac{e_i+sc_i}{b_i}\right\rfloor,
 \qquad K_Y+A_s=\mathcal O_Y(C_s),
 \quad C_s=\sum_Bu_s(B)B.
\]
Only finitely many primes occur with a nonzero slope or intercept.
Replace $d,z$ by a common positive multiple and power so that all
$v_B=\min_i(c_i/b_i)$ are integers, and write
$V=\sum_Bv_BB$.  For all sufficiently large integer $s$,
\begin{equation}\label{inj:affine-stabilization}
 C_s=C_\infty+sV,
 \qquad C_\infty\ge C_0.
\end{equation}
Indeed only the indices of minimal slope can minimize for large
$s$.  Among them the least intercept, after taking its floor,
is at least the unrestricted least intercept at zero.  This proves
the stated inequality coefficient by coefficient.  Set
$A_\infty=\mathcal O_Y(C_\infty)-K_Y$.  Then
$A_s=A_\infty+sV$ eventually, and
\begin{equation}\label{inj:line-inclusion}
 A_0+sV\hookrightarrow A_s
\end{equation}
for those $s$. The eventual family has intercept $A_\infty$, but this
inclusion compares it with the actual zero-degree line $A_0$. It is
this inclusion that will transfer sections produced by interpolation
from exponent zero.

\smallskip\noindent\textit{Integrable metrics and interpolation from zero.}
The original moving part, together with $-M$ pseudoeffective,
shows that $\pi^*L-\epsilon f^*H_1$ is pseudoeffective for some
ample $H_1$ on $Y$ and $\epsilon>0$.  The pullback of $H$ to
$Y$ is big, so it still dominates a small ample class.
Mix a metric representing this pseudoeffectivity with the smooth
semipositive metric on $\pi^*L$, using a small positive coefficient.
This gives at $s=0$ a singular metric with trivial multiplier ideal
and curvature at least $\delta f^*\omega_Y$, with $\delta>0$.
For $s>0$ use the smooth semipositive metrics on
$(1+sd)\pi^*L$.

The singular adjoint metric and its extension
\cite[Theorems 3.3.4--3.3.5]{PT2018} give semipositive metrics
on $\mathscr E_s$ on its locally free locus.  The generic
multiplier-ideal inclusion is an isomorphism since the upstairs
ideal is trivial.  Give $A_s$ the quotient metric for the
holomorphic projection onto this summand; quotient positivity is
\cite[Lemma 2.3.4(2)]{PT2018}.  Write its weights as $\psi_s$.
Their description is the logarithm of the squared norm of the
dual line frame in $\mathscr E_s^\vee$.  They extend
plurisubharmonically across the remaining codimension-two subset
of smooth $Y$, since $A_s$ is a line bundle.  They are not
identically $-\infty$, by the nondegenerate finite fiber metric
at almost every smooth base-change parameter.  An invariant
choice of upstairs metric is unnecessary for this quotient argument.

Crucially,
\begin{equation}\label{inj:base-integrability}
 e^{-\psi_s}\in L^1_{\rm loc}(Y)\qquad(s\ge0).
\end{equation}
A local frame of the actual summand lifts to a local section of
$\mathscr E_s$.  Its fiberwise squared norm bounds the quotient
squared norm from above.  The lift is a holomorphic adjoint section
upstairs, after adding a base canonical frame.  Its absolute
density is integrable over inverse images of relatively compact
charts by the trivial upstairs ideal and properness.  Fubini on
the smooth locus identifies this integral with the fiber norms
integrated against the base canonical volume, proving
\eqref{inj:base-integrability}.  A frame of a reflexive hull
without such a lift would not suffice.

The weight $\psi_0$ has a strictly positive lower curvature bound
on all of $Y$.  To check this across the boundary, choose a smooth
function $p$ on $Y$ whose Hessian dominates $\omega_Y$ near a
given point and is bounded above by $C\omega_Y$ globally.
Multiply the upstairs squared metric by $e^{t f^*p}$ for small
$t>0$.  It remains semipositive and has the same ideal.  Its
direct-image quotient weight is $\psi_0-tp$ almost everywhere,
so it is plurisubharmonic everywhere by the extension theorem.
A finite covering gives $i\partial\bar\partial\psi_0
\ge\delta'\omega_Y$ for some $\delta'>0$.

We now put strictly positive metrics with trivial multiplier ideal
on $A_0+sV$ for every $s\ge0$. These are the lines on the left of
\eqref{inj:line-inclusion}, so vanishing for them will retain the
known section count at zero. Fix an integer $s>0$. Take a sufficiently
large integer $u$ in
the range of \eqref{inj:affine-stabilization}, and choose a smooth
metric $\psi_\infty$ on $A_\infty$.  In compatible frames the
formula
\begin{equation}\label{inj:openness-mixture}
 \rho_s=\psi_0+\frac{s}{u}(\psi_u-\psi_\infty)
\end{equation}
defines a metric on $A_0+sV$.  Its curvature is strictly positive
once $s/u$ is small enough.  Strong openness
\cite[Theorem 1.1]{GuanZhou2015} and compactness give
$\eta>0$ such that $e^{-(1+\eta)\psi_0}$ is locally
integrable on a finite cover.  H\"older, together with
\eqref{inj:base-integrability} for this $u$, proves
$e^{-\rho_s}$ integrable if
\[
 \frac{s}{u}\frac{1+\eta}{\eta}\le1.
\]
The smooth $\psi_\infty$ contributes only bounded factors.
No exponent improvement uniform in $u$ is needed.  At $s=0$
the original $\psi_0$ already has these properties.

Nadel vanishing consequently gives
\[
 H^i(Y,K_Y+A_0+sV)=0\qquad(i>0,\ s\ge0).
\]
The Euler polynomial in $s$ has value one at zero by
\eqref{inj:global-invariant-identification}.  Thus it is nonzero
for all sufficiently large integers $s$.  Using
\eqref{inj:line-inclusion} and then
\eqref{inj:global-invariant-identification} gives a nonzero
invariant section of $sdL$ on $X$.
\end{proof}

Under the hypotheses $H^q(X,\mathcal O_X)=0$ for $q>0$,
Proposition~\ref{har:prop:index} and Lemma~\ref{har:lem:forms} supply the
unbounded invariant twisted forms required to apply this proposition.
For positive differential degree, their saturated generic determinant
has pseudoeffective inverse by the cotangent-tensor theorem, exactly
as in Corollary~\ref{inj:forms-input}, now preserving the action.
For degree zero the invariant section is already available.  Thus
this construction is a separate proof of invariant nonvanishing
under those hypotheses, with strong openness and the actual
direct-image line replacing the bounded base metric used earlier.

Section~\ref{sec:descent} applies the invariant conclusions to compact
factors and descends their products through the retained monodromy
and a finite-cover norm.

\section{Compact monodromy and the descent maps}\label{sec:descent}\label{har:sec:descent}
We prove Corollary~\ref{thm:ordinary} by applying invariant
nonvanishing to the compact factor in the universal cover and then
taking a finite norm. We use the structure statement of \cite[Theorem \textnormal{(Finite cover with compact torus monodromy)}]{CompanionH}, not that companion's nonvanishing conclusion. For a smooth connected projective $X$ with smoothly semipositive $-K_X$, it gives a holomorphic isometric universal-cover splitting
\[
 \widetilde X=\C^c\times V\times F,
\]
where $V$ is a product of compact simply connected Ricci-flat factors with invariant parallel canonical frames and $F$ is a product of compact projective factors with no positive-degree holomorphic forms. After finite index, the deck subgroup acts faithfully by a translation lattice on $\C^c$, trivially on $V$, and through a compact real torus $H$ on $F$. Empty products are points. The associated finite cover $X'$ is smooth projective and need not be a product. These are the only geometric conclusions from the companion that we need.

The compact factor $F$ has $H^j(F,\mathcal O_F)=0$ for $j>0$ by
Hodge symmetry, and its anticanonical metric is smoothly semipositive.
We first descend an invariant section. Afterwards we construct two
additional descent maps, preserving a cohomology class or a twisted
form, to which the ordinary conversions apply on $X'$.

\subsection{Invariant sections and the finite norm}\label{har:subsec:sectiondescent}
\begin{proof}[Proof of Corollary~\ref{thm:ordinary}]
On $F$, the compact torus preserves a very ample line because it preserves its Chern class and $\operatorname{Pic}^0(F)=0$. Its projective action has algebraic-torus Zariski closure: a compact connected abelian subgroup of a projective linear group is conjugate into a diagonal compact torus. Passing through the finite special-linear cover, if necessary, provides a linearization of a very ample power. Invariants of the natural anticanonical line for the compact group and its Zariski closure agree. Corollary~\ref{thm:main} therefore gives a nonzero invariant section of $-mK_F$ for some $m>0$.

Tensor its pullback to $\widetilde X$ with the $m$-th powers of the inverse canonical frames on $\C^c$ and $V$. This nonzero product is invariant under the selected deck subgroup and descends to a section on $X'$. If $F$ is a point, take its section to be $1$ and $m=1$.

The finite \'etale norm \cite[Lemma \textnormal{(Finite \'etale norm)}]{CompanionH} sends this section to a nonzero section of a positive power of $-K_X$. Explicitly, pull to a connected finite Galois refinement, multiply all Galois translates, and descend their invariant product. The product is nonzero on the integral cover, and canonical bundles pull back under \'etale morphisms. Projectivity and GAGA make the resulting holomorphic section algebraic.
\end{proof}

The exponent can increase when common powers are taken on compact factors and when the finite norm is formed. None of these arguments supplies a uniform exponent.

\subsection{Descending a nonzero cohomology class}\label{har:subsec:cohomdescent}

We can also descend cohomology before producing any section.
The invariant Euler characteristic of $F$ is one. Proposition~\ref{har:prop:index} gives, for unbounded \(m\),
nonzero classes in \(H^q(F,-mK_F)^H\), with one fixed \(q\) after
subsequence extraction. Choose a smooth closed Dolbeault representative
and average it over \(H\). Pull it to the product universal cover and
tensor with powers of the constant anticanonical volume on
\(\mathbb C^c\) and the invariant trivializations on \(V\). It is
deck-invariant for the selected subgroup and descends to the finite
cover \(X'\).

The descended class is nonzero. The faithful translation action makes
a compact slice \(F\), at fixed Euclidean and \(V\) coordinates,
embed into \(X'\); no nonidentity lattice element preserves that
slice. Restricting the descended coefficient-valued Dolbeault class
to it, with the product coefficient frames fixed, recovers the chosen
nonzero class on \(F\). Hence \(H^q(X',-mK_{X'})\ne0\) for unbounded
\(m\). Theorem~\ref{thm:alteration-conversion} now produces a positive
anticanonical section on \(X'\). The finite norm of
Section~\ref{har:subsec:sectiondescent} gives a positive anticanonical
section on $X$. Here the invariant cohomology is descended before the
ordinary alteration conversion is applied on $X'$.

\subsection{Descending twisted forms before Iitaka conversion}
\label{har:subsec:formdescent}

There is a direct twisted-form route to the ordinary conversion of
Section~\ref{inj:injectivity}.  Keep the product
$\widetilde X=\mathbb C^c\times V\times F$ and the finite-index
deck subgroup just constructed, whose action is by translations on
$\mathbb C^c$, trivial on $V$, and contained in the compact torus
$H$ on $F$.  Put $P=-K_F$ and $d=\dim F$.  The invariant index has
constant term one.  Hence for a fixed $q$ and unbounded positive
integers $m$ there are nonzero invariant classes in
\[
 H^q(F,mP)^H=H^q(F,K_F+(m+1)P)^H.
\]
Choose an $H$-invariant K\"ahler form.  Smooth semipositive hard
Lefschetz \cite{DPS2001} is an equivariant surjection from
$H^0(F,\Omega_F^{d-q}\otimes (m+1)P)$ onto the displayed group.
Averaging a preimage over $H$ keeps its nonzero image, so supplies
\[
 0\ne\alpha_m\in
 H^0(F,\Omega_F^{d-q}\otimes (m+1)P)^H.
\]
Set $p=d-q$.  Let $\tau_0$ be the translation-invariant inverse
canonical frame on $\mathbb C^c$, and let $\tau_V$ be a product of
the parallel inverse canonical frames on $V$.  Projection to $F$
and the natural inclusion of its pulled-back cotangent bundle into
the product cotangent bundle give the holomorphic twisted form
\[
 \widetilde\alpha_m=
 \operatorname{pr}_F^*\alpha_m\otimes
 (\tau_0\otimes\tau_V)^{m+1}
 \in H^0\bigl(\widetilde X,
       \Omega_{\widetilde X}^{p}\otimes(-K_{\widetilde X})^{m+1}\bigr).
\]
It is nonzero: evaluate $\alpha_m$ on a tangent $p$-tuple at a point
where it is nonzero and use the same tuple in the product; the other
coefficient frames never vanish.  The $m+1$ coefficient exponent
matches the full anticanonical line of the product.

Every factor in this expression is invariant under the selected
deck subgroup.  The resulting form therefore descends through the
covering map to a holomorphic section
\[
 0\ne\beta_m\in
 H^0\bigl(X',\Omega_{X'}^{p}\otimes(-K_{X'})^{m+1}\bigr)
\]
on the associated smooth projective finite \'etale cover $X'$ of
$X$.  Descent follows on evenly covered coordinate neighborhoods,
where invariant local forms glue; its pullback is
$\widetilde\alpha_m$, so the descended form cannot be zero.
For every fixed $r>0$ the same construction is linear and gives
the injection
\[
 H^0(F,\Omega_F^p\otimes rP)^H
 \lhook\joinrel\longrightarrow
 H^0(X',\Omega_{X'}^p\otimes(-rK_{X'})).
\]
Projectivity and GAGA make $\beta_m$ algebraic.  If $F$ is a point,
take $p=q=0$ and $\alpha_m=1$; the same construction applies.

Now apply Corollary~\ref{inj:forms-input} on $X'$ itself to this
fixed differential degree and the unbounded positive twists $m+1$.
It gives a nonzero positive anticanonical section on $X'$.  Pull back
to a finite connected Galois refinement and multiply its translates,
as in Section~\ref{har:subsec:sectiondescent}, to obtain a nonzero
positive anticanonical section on $X$.  Thus this route descends the
twisted forms before using the ordinary Iitaka and injectivity
conversion; the cohomology and alteration route of the preceding
subsection remains a separate construction.

\subsection{The stronger tensor-holonomy alternative}
The Bochner--holonomy argument underlies the structure results of
\cite[Sections~2--3]{DPS1996} and
\cite[proof of Theorem~1.4]{CDP2015}. We include the following tensor
form of that argument as a second input for the same descent. It
excludes all positive covariant tensor powers on an irreducible
non-Ricci-flat factor, which is stronger than excluding holomorphic
differential forms. The proof uses the determinant character of the
irreducible holonomy representation.

\begin{proposition}[The covariant-tensor holonomy dichotomy]
\label{prop:tensor-holonomy}
Let $P$ be a compact simply connected nonflat K\"ahler factor,
irreducible in the Riemannian de Rham sense, with nonnegative Ricci
curvature.  Either $K_P$ has a nonzero
parallel holomorphic volume form, or
\[
 H^0\bigl(P,(\Omega_P^1)^{\otimes m}\bigr)=0
 \qquad\text{for every integer }m>0.
\]
\end{proposition}

\begin{proof}
For a holomorphic covariant tensor, the contracted curvature on its
bundle is the sum, over its tensor slots, of the negative Ricci
operator.  It is nonpositive.  Integration of the Bochner identity
on compact $P$ therefore makes the tensor parallel.  Its value at
a point must consequently be fixed by the holonomy representation.

Let $H$ be the closure of the holonomy group on the complex tangent
space at that point.  It is a compact subgroup of the unitary group.
Simple connectivity makes full holonomy equal to restricted holonomy,
so $H$ is connected.  De Rham irreducibility makes its complex tangent
representation irreducible: a proper invariant complex subspace
would give a proper real parallel subspace and its orthogonal
complement.  Consider the determinant character of this
representation.  If it is trivial, the induced holonomy on the
canonical line is trivial.  Parallel transport of a nonzero canonical
vector gives a global parallel volume form, which is holomorphic.

If the determinant character is nontrivial, connectedness implies
that its differential is nonzero.  This differential is the complex
trace on the Lie algebra $\mathfrak h$.  A compact Lie algebra is the
sum of its center and its derived algebra, and trace vanishes on the
latter.  Hence there is a central $Z\in\mathfrak h$ with nonzero
trace.  By complex irreducibility and Schur's lemma it acts on the
tangent space as a scalar $ia\,\mathrm{Id}$, with $a\in\mathbb R$
nonzero.  Its action on the $m$-fold covariant tensor power is the
nonzero scalar $-mia$.  That tensor power has no $H$-fixed vector.
The parallelism proved in the first paragraph now excludes every
nonzero holomorphic tensor of positive degree.
\end{proof}

This dichotomy supplies a second factor argument, stronger than
vanishing of differential forms.  On its trivial-canonical factors,
a connected group of holomorphic isometries fixes a holomorphic
volume form: it acts trivially on cohomology by isotopy, and the
holomorphic top-form space injects into cohomology by Hodge theory.
On the other factors all positive-degree forms vanish as direct
summands of covariant tensor powers.  Hodge symmetry then gives
vanishing of higher structure-sheaf cohomology.  The compact-monodromy
algebraization, invariant section, and norm construction above apply
to this dichotomy as well.

\clearpage
\bibliographystyle{amsplain}
\bibliography{references}
\end{document}
